\exercisesection{}

除另有说明外，在本节的各习题中，E 表示一个复希尔伯特空间。

\begin{enumerate}
\def\labelenumi{\arabic{enumi})}
\tightlist
\item
  设 \(u\) 是 E 的自同态，其像在 E 中是闭的。设 \(q\) 是 E 的正交投影算子，并设 \(v\) 是 E 的像的自同态 \(x \mapsto q(u(x))\) 。
\end{enumerate}

\begin{enumerate}
\def\labelenumi{\alph{enumi})}
\item
  \(v\) 的伴随是 E 的像的自同态 \(x \mapsto q(u^{*}(x))\) 。
\item
  若 \(u\) 是埃尔米特的，则 \(v\) 是埃尔米特的。
\item
  有可能 \(u\) 是正规的而 \(v\) 不是正规的。
\end{enumerate}

\begin{enumerate}
\def\labelenumi{\arabic{enumi})}
\setcounter{enumi}{1}
\item
  直接利用谱定理（IV, p. 193 的定理 1 的推论）证明 IV, p. 149 的定理 1。
\item
  设 \(S \subset \mathbf{C}\) 是 \(\mathbf{C}\) 的一个紧子集。设 \(\varphi \colon \mathcal{C}(S) \to \mathcal{L}(E)\) 是对合代数的一个连续的含单位元态射。存在唯一的正规自同态 \(u\) ，它作用在 \(E\) 上，谱为 \(S\)，并且使得 \(\varphi(f) = f(u)\) 对所有 \(f \in \mathcal{L}^{\infty}(S)\) 成立。
\item
  设 \(u \in \mathcal{L}(\mathrm{E})\) 是正规的。\(u\) 的谱的凸包是 \(\mathbf{C}\) 中数值值域 \(\iota(u)\) （相应于 \(u\) ；参见 I, p. 135 的定义 2）的闭包。
\item
  设 \(u \in \mathcal{L}(\mathrm{E})\) 满足 \(\| u \| \leqslant 1\) 。对每个 \(\lambda \in \mathbf{C}\) ，用 \(p_{\lambda}\) 表示 E 到 \(u\) 的对应于 \(\lambda\) 的特征空间上的正交投影算子。对 E 中所有的 \(x\) 与 \(y\)，我们有
\end{enumerate}

\[
\lim_{N\to +\infty}\frac{1}{N}\sum_{n = 1}^{N}|\langle x \mid u^{n}(y)\rangle |^{2} = \sum_{\substack{\lambda \in \mathbf{C}\\ |\lambda | = 1}}|\langle p_{\lambda}(y) \mid p_{\lambda}(x)\rangle^{2}|\\ = \sum_{\substack{\lambda \in \mathbf{C}\\ |\lambda | = 1}}|\langle y \mid p_{\lambda}(x)\rangle |^{2}
\]

（可以利用 II, p. 300 的习题 64）。

\begin{enumerate}
\def\labelenumi{\arabic{enumi})}
\setcounter{enumi}{5}
\tightlist
\item
  设 E 是预希尔伯特空间，F 是巴拿赫空间，\(u: E \rightarrow F\) 是线性映射。
\end{enumerate}

\begin{enumerate}
\def\labelenumi{\alph{enumi})}
\item
  证明：u 是紧线性映射当且仅当 \(\|u(e_{n})\|\to0\) 对 E 的每个标准正交序列 \((e_{n})\) 成立。（为证明该条件是充分的，利用 TG, IX, p. 20 的命题 14。）
\item
  假设 E 是可分希尔伯特空间且 F = E。为使存在 E 的标准正交基 \((e_{n})\) 使得 \(\|u(e_{n})\| \to 0\)，必须且只需 u 不是弗雷德霍姆算子。
\item
  假设 E 是希尔伯特空间，并且 \(\langle u(e_{n}), e_{n} \rangle \to 0\) 对 H 的每个标准正交序列 \((e_{n})\) 成立。证明 u 是紧线性映射。
\end{enumerate}

\begin{enumerate}
\def\labelenumi{\arabic{enumi})}
\setcounter{enumi}{6}
\tightlist
\item
  设 \(u\) 与 \(v\) 是 E 的连续自同态。
\end{enumerate}

\begin{enumerate}
\def\labelenumi{\alph{enumi})}
\item
  证明 \(\operatorname{Im}(u) + \operatorname{Im}(v) = \operatorname{Im}((u \circ u^* + v \circ v^*)^{1/2})\) 。
\item
  若 \(u\) 是正的且具有闭像，证明 \(\operatorname{Im}(u) = \operatorname{Im}(u^{1/2})\) 。
\item
  若 \(u, v\) 是正的且具有闭像，并且 \(u + v\) 也具有闭像，则我们有 \(\operatorname{Im}(u + v) = \operatorname{Im}(u) + \operatorname{Im}(v)\) 。
\end{enumerate}

\begin{enumerate}
\def\labelenumi{\arabic{enumi})}
\setcounter{enumi}{7}
\tightlist
\item
  设 F 是复希尔伯特空间，\(u \colon \mathrm{E} \to \mathrm{F}\) 是具有闭像的连续线性映射。
\end{enumerate}

\begin{enumerate}
\def\labelenumi{\alph{enumi})}
\tightlist
\item
  证明存在唯一的连续线性映射 \(v \in \mathcal{L}(\mathrm{F};\mathrm{E})\) 满足下列条件：
\end{enumerate}

\begin{enumerate}
\def\labelenumi{(\roman{enumi})}
\item
  \(u \circ v \circ u = u\) 且 \(v \circ u \circ v = v\) ；
\item
  \(u \circ v\) 与 \(v \circ u\) 是自伴的。
\end{enumerate}

我们把这一映射记作 \(u^{\dagger}\) （穆尔–彭罗斯伪逆）。

\begin{enumerate}
\def\labelenumi{\alph{enumi})}
\setcounter{enumi}{1}
\tightlist
\item
  证明 \(\operatorname{Im}(u) = \operatorname{Im}(u \circ u^{\dagger})\)。证明：对 \(x \in E\) 与 \(y \in \operatorname{Im}(u)\) ，条件 \(u(x) = y\) 与 \(x - u^{\dagger}(y) \in \operatorname{Im}(1_{\operatorname{E}} - u^{\dagger} \circ u)\) 等价。
\end{enumerate}

\begin{enumerate}
\def\labelenumi{\arabic{enumi})}
\setcounter{enumi}{8}
\tightlist
\item
  设 \(u\) 与 \(v\) 是 E 的具有闭像的自同态。若 \(u + v\) 也具有闭像，我们令 \(u: v = u \circ (u + v)^{\dagger} \circ v\) ，其中 \((u + v)^{\dagger}\) 是 \(u + v\) 的穆尔–彭罗斯伪逆。
\end{enumerate}

\begin{enumerate}
\def\labelenumi{\alph{enumi})}
\item
  证明 \(u:v\) 是自伴的。
\item
  我们有 \(u:v=v:u\) 且 \(\operatorname{Im}(u:v)=\operatorname{Im}(u)\cap\operatorname{Im}(v)\) 。
\item
  假设 u 与 v 是正交投影算子。证明 2u : v 是以 \(\operatorname{Im}(u) \cap \operatorname{Im}(v)\) 为像的正交投影算子。
\end{enumerate}

\begin{enumerate}
\def\labelenumi{\arabic{enumi})}
\setcounter{enumi}{9}
\tightlist
\item
  E 的连续自同态 \(u\) 称为亚正规的，如果 \(u^{*} \circ u \geqslant u \circ u^{*}\)；称为仿正规的，如果 \(\|u^{2}(x)\|\|x\| \geqslant \|u(x)\|^{2}\) 对所有 \(x \in E\) 成立。
\end{enumerate}

\begin{enumerate}
\def\labelenumi{\alph{enumi})}
\item
  若 \(u\) 是仿正规的，则 \(u^n\) 对每个 \(n \in \mathbf{N}\) 都是仿正规的。
\item
  若 u 是仿正规的且可逆，则 \(u^{-1}\) 是仿正规的。
\item
  亚正规自同态是仿正规的。
\item
  设 u 是谱包含在单位圆周中的仿正规自同态；则 u 是酉的。
\item
  若 \(u\) 是仿正规自同态，则它的范数 \(\| u\|\) 等于它的谱半径 \(\varrho(u)\) 。
\item
  我们假设 u 是仿正规的。
\end{enumerate}

\begin{enumerate}
\def\labelenumi{\arabic{enumi})}
\setcounter{enumi}{10}
\tightlist
\item
  设 u 是 E 的连续自同态。
\end{enumerate}

\begin{enumerate}
\def\labelenumi{\alph{enumi})}
\tightlist
\item
  为使 \(u\) 是仿正规的，必须且只需
\end{enumerate}

\[
(u ^ {*}) ^ {2} \circ u ^ {2} - 2 t (u ^ {*} \circ u) + t ^ {2} 1 _ {\mathrm{E}} \geqslant 0
\]

对每个实数 t \textgreater{} 0 成立。

\begin{enumerate}
\def\labelenumi{\alph{enumi})}
\setcounter{enumi}{1}
\tightlist
\item
  为使 \(u\) 是仿正规的，必须且只需
\end{enumerate}

\[
p \circ q ^ {2} \circ p - 2 t p ^ {2} + t 1 _ {\mathrm{E}} \geqslant 0
\]

对所有 \(t > 0\) 成立，其中 \(p = (u \circ u^{*})^{1/2}\) ，\(q = (u^{*} \circ u)^{1/2}\) 。

\begin{enumerate}
\def\labelenumi{\arabic{enumi})}
\setcounter{enumi}{11}
\tightlist
\item
  \begin{enumerate}
  \def\labelenumii{\alph{enumii})}
  \tightlist
  \item
    设 u 与 v 是 E 的连续正自同态。假设对每个 t \textgreater{} 0，我们有
  \end{enumerate}
\end{enumerate}

\[
2 t u ^ {2} \circ (u ^ {2} + t ^ {2} 1 _ {\mathrm{E}}) ^ {- 1} \leqslant v \leqslant \frac {1}{2 t} (u ^ {2} + t ^ {2} 1 _ {\mathrm{E}}).
\]

证明 u = v。

\begin{enumerate}
\def\labelenumi{\alph{enumi})}
\setcounter{enumi}{1}
\tightlist
\item
  为使 \(u\) 是正规自同态，必须且只需 \(u\) 与 \(u^*\) 都是仿正规的并且有相同的核。
\end{enumerate}

\begin{enumerate}
\def\labelenumi{\arabic{enumi})}
\setcounter{enumi}{12}
\tightlist
\item
  \begin{enumerate}
  \def\labelenumii{\alph{enumii})}
  \tightlist
  \item
    设 A 是 E 的单位向量所成的无限集，它沿有限子集的余集滤子弱收敛于 0。证明存在取值于 A 的序列 \((x_{n})_{n\geqslant 1}\) 以及 E 中的标准正交序列 \((e_n)_{n\geqslant 1}\) ，使得 \(\| x_n - e_n\| \leqslant 1 / n\) 对所有 \(n\geqslant 1\) 成立 。
  \end{enumerate}
\end{enumerate}

\begin{enumerate}
\def\labelenumi{\alph{enumi})}
\setcounter{enumi}{1}
\item
  设 u 是 E 的自同态，未必连续。假设对 E 中每个标准正交序列 \((e_{n})\) ，序列 \((u(e_{n}))\) 都弱收敛于 0。证明 u 是连续且紧的。
\item
  设 u 是 E 的连续自同态。为使存在 E 的标准正交基 \((e_{n})\) 使序列 \((u(e_{n}))\) 弱收敛于 0，必须且只需：对 E 的每个自同态 v，\(v \circ u - 1_{E}\) 都不是紧的。
\end{enumerate}

\begin{enumerate}
\def\labelenumi{\arabic{enumi})}
\setcounter{enumi}{13}
\tightlist
\item
  设 X 是局部紧空间，\(\mu\) 是 X 上的正有界测度且 \(\mu(\mathrm{X}) = 1\) ，并设 \(f: X \to X\) 是 \(\mu\)-可测映射且满足 \(f(\mu) = \mu\) 。我们用 \(u_{f}\) 表示 \(\mathrm{L}_{\mathbf{C}}^{2}(\mathrm{X}, \mu)\) 的由 \(\varphi \mapsto f \circ \varphi\) 定义的酉自同态（参见 I, p. 190 的习题 16）。对每个 \(n \in N\)，我们记 \(f^{n} = f^{\circ n}\) （\(n\) 次复合），并且对 X 的每个子集 A，记 \(f^{-n}(\mathrm{A}) = g^{-1}(\mathrm{A})\) ，其中 \(g = f^{n}\) 。
\end{enumerate}

我们称 \(f\) 关于 \(\mu\) 是弱混合的，如果对 X 的所有 \(\mu\)-可测子集 A 与 B，我们有

\[
\lim _ {N \rightarrow + \infty} \frac {1}{N} \sum_ {n = 0} ^ {N} | \mu (A \cap f ^ {- n} (B) - \mu (A) \mu (B) | = 0.
\]

\begin{enumerate}
\def\labelenumi{\alph{enumi})}
\tightlist
\item
  下列条件等价：
\end{enumerate}

\begin{enumerate}
\def\labelenumi{(\roman{enumi})}
\item
  映射 f 关于 \(\mu\) 是弱混合的；
\item
  对所有 \(\varphi_{1},\varphi_{2} \in \mathrm{L}_{\mathbf{C}}^{2}(\mathrm{X}, \mu)\) ，我们有
\end{enumerate}

\[
\lim _ {N \to + \infty} \frac {1}{N} \sum_ {n = 0} ^ {N} | \langle \varphi_ {1} | u _ {f} ^ {n} (\varphi_ {2}) \rangle - \langle 1 | \varphi_ {2} \rangle \langle \varphi_ {1} | 1 \rangle | = 0;
\]

\begin{enumerate}
\def\labelenumi{(\roman{enumi})}
\setcounter{enumi}{2}
\item
  映射 \(f \times f: X \times X \to X \times X\) 是 \((\mu \otimes \mu)\) -遍历的（参见 I, p. 190 的习题 16）；
\item
  相应于特征值 1 的特征空间的维数为 1，且对每个 \(\lambda \in \mathbf{C} - \{1\}\) ，\(u_{f}\) 相应于 \(\lambda\) 的特征空间化为 0。
\end{enumerate}

（为证明 (iii) 蕴含 (i)，利用 II, p. 300 的习题 64。）

\begin{enumerate}
\def\labelenumi{\alph{enumi})}
\setcounter{enumi}{1}
\item
  考虑 I, p. 190 习题 16 中 e) 部分的例子。映射 f 关于 \(\mu\) 是弱混合的。
\item
  考虑 I, p. 190 习题 16 中问题 f) 的例子。映射 f 关于 \(\mu\) 不是弱混合的。
\end{enumerate}

\begin{enumerate}
\def\labelenumi{\arabic{enumi})}
\setcounter{enumi}{14}
\tightlist
\item
  设 \(G = (S, F, o, t, \iota)\) 是一个图（TA, II, p. 155 的定义 1）。我们假设 G 是连通的，S 是可数的，并且 S 至少有两个元素。对每个 \(x \in S\) ，我们用 \(v(x)\) 表示 G 的以 x 为端点的边的数目；我们有 \(v(x) \geqslant 1\) 对所有 \(x \in S\) 成立，并且假设 \(v(x)\) 对所有 \(x \in S\) 是有限的。对 S 中所有 x 与 y，我们用 \(a(x, y)\) 表示 G 中以 x 与 y 为端点的箭的数目。我们假设 \(a(x, y)\) 对所有 \((x, y) \in S^{2}\) 是有限的。
\end{enumerate}

我们定义 S 上的测度 \(\mu\) 如下：\(\mu(\{x\}) = v(x)\) 对所有 \(x \in S\) 成立。测度 \(\mu\) 是正的。我们用 \(\mathcal{F}(\mathrm{G})\) 表示 S 上取复值的函数所成的向量空间。我们用 \(\widetilde{\mathrm{M}}_{\mathrm{G}}\) 表示 \(\mathcal{F}(\mathrm{G})\) 的满足下式的自同态

\[
\widetilde {\mathrm{M}} _ {\mathrm{G}} (\varphi) (x) = \frac {1}{v (x)} \sum_ {y} a (x, y) \varphi (y),
\]

对所有 \(\varphi\in\mathcal{F}(\mathrm{G})\) 及每个 \(x\in\mathrm{S}\) 成立（“马尔可夫算子”）。

\begin{enumerate}
\def\labelenumi{\alph{enumi})}
\item
  自同态 \(\widetilde{\mathrm{M}}_{\mathrm{G}}\) 通过向子空间的过渡诱导出 \(\mathrm{M_G}\) ，即 \(\mathrm{L}_{\mathbf{C}}^{2}(\mathrm{S},\mu)\) 的一个自同态。
\item
  自同态 \(M_{G}\) 是自伴的。它满足
\end{enumerate}

\[
\langle \varphi | 1 - \mathrm{M} _ {\mathrm{G}} (\varphi) \rangle = \frac {1}{2} \sum_ {(x, y) \in \mathrm{S} ^ {2}} a (x, y) | \varphi (x) - \varphi (y) | ^ {2}
\]

\[
\langle \varphi | 1 + \mathrm{M} _ {\mathrm{G}} (\varphi) \rangle = \frac {1}{2} \sum_ {(x, y) \in \mathrm{S} ^ {2}} a (x, y) | \varphi (x) + \varphi (y) | ^ {2}
\]

对所有 \(\varphi\in\mathrm{L}^{2}(\mathrm{S},\mu)\) 成立。此外，我们有 \(-1\leqslant\mathrm{M}_{\mathrm{G}}\leqslant1\) 。

\begin{enumerate}
\def\labelenumi{\alph{enumi})}
\setcounter{enumi}{2}
\tightlist
\item
  设 \(d \geqslant 1\) 是整数，并假设 G 是 \(Z^{d}\) 相对于集合 \(\{e_{i}\}\) 的凯莱图，其中 \((e_{i})\) 是 \(Z^{d}\) 的典范基（TA, II, p. 221 的习题 7）。\(\varrho(\mathrm{M}_{\mathrm{G}})\) （\(M_{G}\) 的谱半径）等于 1。
\end{enumerate}

\(d\) ) 设 \(d \geqslant 3\) 是整数。假设 G 是一棵 \(d\) -正则的无限树（TA, II, p. 157），即其中每个顶点都有 \(d\) 个邻点的无限树。设 \(x_0\) 是 G 的一个固定顶点。对每个 \(\lambda \in \mathbf{C}\) ，存在唯一的函数 \(\varphi_{\lambda} \colon S \to \mathbf{C}\) ，使得 \(\varphi_{\lambda}(x_0) = 1\) 且 \(\widetilde{\mathrm{M}}_{\mathrm{G}}(\varphi_{\lambda}) = \lambda \varphi_{\lambda}\) 。此外，对每个 \(x \in S\) ，\(\varphi_{\lambda}\) 在 \(x\) 处的值只依赖于 \(x\) 到 \(x_0\) 的距离（参见 TA, II, p. 222 的习题 10）。

\begin{enumerate}
\def\labelenumi{\alph{enumi})}
\setcounter{enumi}{4}
\tightlist
\item
  \(M_{G}\) 的谱是区间 \([-2\sqrt{d-1}/d, 2\sqrt{d-1}/d]\) （“凯斯滕定理”）。特别地，我们有 \(\varrho(\mathrm{M}_{\mathrm{G}})=2\sqrt{d-1}/d\) 。
\end{enumerate}

\begin{enumerate}
\def\labelenumi{\arabic{enumi})}
\setcounter{enumi}{15}
\tightlist
\item
  我们沿用上一个习题的记号，并假设 S 与 F 是有限的。我们用 \(h_{G}\) 表示图 G 的切格常数，它定义为
\end{enumerate}

\[
h _ {\mathrm{G}} = \inf _ {\varnothing \neq \mathrm{V} \subset \mathrm{S}} \frac {\operatorname{Card} (\mathrm{E} (\mathrm{V} , \mathrm{S} - \mathrm{V}))}{\inf (\operatorname{Card} (\mathrm{V}) , \operatorname{Card} (\mathrm{S} - \mathrm{V}))},
\]

其中 E(V, W) 表示 G 的一个端点在 V 中、另一个端点在 W 中的箭所成的集。

\begin{enumerate}
\def\labelenumi{\alph{enumi})}
\tightlist
\item
  我们有 \(\varrho(\mathrm{M}_{\mathrm{G}})=1\) ，并且 \(\lambda=1\) 是 \(\mathrm{M}_{\mathrm{G}}\) 的重数为 1 的特征值。我们有 \(-1\in\mathrm{Sp}(\mathrm{M}_{\mathrm{G}})\) 当且仅当 G 是二部的（参见 TA, II, p. 219 的习题 2）；这时，\(\mathrm{M}_{\mathrm{G}}\) 的谱在 \(\lambda\mapsto-\lambda\) 下不变，并且 \(\lambda\) 的谱重数等于 \(-\lambda\) 的谱重数对所有 \(\lambda\in\mathbf{C}\) 成立。
\end{enumerate}

我们用 \(\mathrm{L}_{0}^{2}(\mathrm{S},\nu)\) 表示在 \(\mathrm{L}_{\mathbf{C}}^{2}(\mathrm{S},\nu)\) 中、\(M_{G}\) 的相应于特征值 1 与 -1 的诸特征空间之并的正交补。通过向子空间的过渡诱导的自同态 \(M_{G}\) 是 \(\mathrm{L}_{0}^{2}(\mathrm{S},\nu)\) 的自同态，记作 \(M_{G,0}\)。我们有 \(\varrho(\mathrm{M}_{\mathrm{G},0}) < 1\)。我们用 \(\lambda_{1}(\mathrm{G})\) 表示 \(1 - M_{G}\) 的最小非零特征值。

\begin{enumerate}
\def\labelenumi{\alph{enumi})}
\setcounter{enumi}{1}
\tightlist
\item
  设 \(v_{+} = \sup v(x)\) ，\(v_{-} = \inf v(x)\) 。我们有
\end{enumerate}

\[
\lambda_ {1} (\mathrm{G}) \leqslant \frac {2 v _ {+}}{v _ {-} ^ {2}} h _ {\mathrm{G}}.
\]

\begin{enumerate}
\def\labelenumi{\alph{enumi})}
\setcounter{enumi}{2}
\tightlist
\item
  设 \(\varphi\colon S\to\mathbf{R}\) 是 \(S\) 上的非常数实值函数；令 \(m=\inf\varphi(x)\) ，\(M=\sup\varphi(x)\) 。设 \(t_{0}\in\mathbf{R}\) 使得：对每个 \(t\in\mathbf{R}\) ，集合 \(W_{t}=\overline{\varphi}^{-1}(]-\infty,t_{0}[)\) 满足 \(\mathrm{Card}(W_{t})\leqslant\mathrm{Card}(S)/2\) 当且仅当 \(t<t_{0}\) 。设 \(\nu\) 是支撑含于 \([m,M]\) 的无原子测度（INT, V, p. 67 的定义 5）。令
\end{enumerate}

\[
\mathrm{A} = \frac {1}{2} \sum_ {(x, y) \in \mathrm{S} ^ {2}} a (x, y) \nu ([ \varphi (x), \varphi (y) ])
\]

\[
\mathrm{B} = \sum_ {x \in \mathrm{S}} \nu ([ t _ {0}, \varphi (x) ]),
\]

其中，当 \(a < b\) 是实数时，我们记 \(\nu([b, a]) = \nu([a, b])\) 。存在 \(t \in \mathbf{R}\) 使得

\[
\frac {\operatorname{Card} \left(\mathrm{E} \left(\mathrm{W} _ {t} , \mathrm{S} - \mathrm{W} _ {t}\right) \right.}{\inf \left(\operatorname{Card} \left(\mathrm{W} _ {t}\right) , \operatorname{Card} \left(\mathrm{S} - \mathrm{W} _ {t}\right)\right)} \leqslant \frac {\mathrm{A}}{\mathrm{B}}.
\]

\begin{enumerate}
\def\labelenumi{\alph{enumi})}
\setcounter{enumi}{3}
\tightlist
\item
  我们有
\end{enumerate}

\[
h _ {\mathrm{G}} \leqslant v _ {+} \sqrt {2 \lambda_ {1} (\mathrm{G})}.
\]

（把上一问题应用于 \(\varphi\) 取为 \(\mathrm{M_G}\) 的相应于特征值 \(\lambda_1(\mathrm{G})\) 的特征函数并且 \(\nu\) 取适当选取的测度的情形。）

\begin{enumerate}
\def\labelenumi{\alph{enumi})}
\setcounter{enumi}{4}
\tightlist
\item
  我们称一个族 \((\mathrm{G}_{j})_{j\in\mathrm{J}}\) （由有限图 \(\mathrm{G}_{j}=(\mathrm{S}_{j},\mathrm{F}_{j},o_{j},t_{j},i_{j})\) 组成）为扩张图族，如果下列条件成立：
\end{enumerate}

\begin{enumerate}
\def\labelenumi{(\roman{enumi})}
\item
  沿着 J 的有限子集的余集滤子，我们有 \(\mathrm{Card}(\mathrm{S}_{j}) \to +\infty\) ；
\item
  存在 \(C \geqslant 0\) ，使得对每个 j 及每个 \(x \in S_{j}\) ，\(G_{j}\) 在 x 处的度 \(\leqslant C\) ；
\item
  存在 \(\delta > 0\) ，使得 \(h_{\mathrm{G}_j} \geqslant \delta\) 对所有 \(j \in \mathrm{J}\) 成立。
\end{enumerate}

证明 \((\mathrm{G}_{j})\) 是扩张图族当且仅当 (i)、(ii) 以及条件

(iii') 存在 \(\delta > 0\) ，使得 \(\lambda_1(G_j) \geqslant \delta\) 对所有 \(j \in J\) 成立。

\begin{enumerate}
\def\labelenumi{\arabic{enumi})}
\setcounter{enumi}{16}
\tightlist
\item
  \begin{enumerate}
  \def\labelenumii{\alph{enumii})}
  \tightlist
  \item
    设 \(n \in \mathbf{N}\)。设 E 是维数为 \(n\) 的希尔伯特空间，\(u \in \mathcal{L}(\mathrm{E})\) ，并设 \(x \in \mathrm{E}\) 是范数为 1 的向量且满足 \(u(x) = 0\)。对每个范数为 1 的向量 \(y \in \mathrm{E}\) ，我们用 \(\Delta_u(y)\) 表示在 \(y^\circ\) 上由 \(z \mapsto \langle z | u(z) \rangle\) 定义的二次形式的判别式。
  \end{enumerate}
\end{enumerate}

对每个范数为 1 的 \(y \in E\) ，我们有 \(|\langle x | y \rangle|^{2} \Delta_{u}(x) = \Delta_{u}(y)\) 。（注意，\(\Delta_{u}(y) = \langle a | \bigwedge u^{n-1}(a) \rangle / \langle a | a \rangle\) 对空间 \(\bigwedge^{n-1} y^{\circ} \subset \bigwedge^{n-1} E\) 的每个非零元素 a 成立，其中空间 \(\bigwedge^{n-1} E\) 赋有其典范希尔伯特结构，参见 EVT, V, p.33；然后分别在 x 与 y 正交以及 x = y 的情形验证该公式。）

\begin{enumerate}
\def\labelenumi{\alph{enumi})}
\setcounter{enumi}{1}
\tightlist
\item
  设 A 是复系数的 \((n,n)\) 型矩阵，满足 \({}^{t}\overline{A}=A\) 。对每个满足 \(1\leqslant i\leqslant n\) 的整数 i，用 \(A_{i}\) 表示 \((n-1,n-1)\) 型的、通过删除 A 的第 \(i\) 行与第 \(i\) 列所得的矩阵。设 \((e_{1},\ldots,e_{n})\) 是 \(C^{n}\) 的、赋有其典范希尔伯特结构且由 A 的特征向量组成的标准正交基，并用 \(\lambda_{i}\) 表示 u 相应于 \(e_{i}\) 的特征值。记 \(e_{i}=(e_{i,j})_{1\leqslant j\leqslant n}\) 。
\end{enumerate}

对所有 \((i,j)\) ，我们有

\[
\left| e _ {i, j} \right| ^ {2} \prod_ {k \neq i} \left(\lambda_ {i} - \lambda_ {k}\right) = \det \left(\lambda_ {i} - A _ {j}\right).
\]

（化归到 \(\lambda_{i}=0\) 的情形，然后把上一问题应用于空间 \(C^{n}\) 、自同态 \(u\colon y\mapsto Ay\) 以及向量 \(x=e_{i}\) 。）

\begin{enumerate}
\def\labelenumi{\arabic{enumi})}
\setcounter{enumi}{17}
\tightlist
\item
  设 X 是紧可度量化空间，\(\mu\) 是 X 上的正测度。记 \(\mathrm{E} = \mathrm{L}_{\mathbf{C}}^{2}(\mathrm{X}, \mu)\)。设 \((\mathcal{X}_{n})_{n \in \mathbf{N}}\) 是把 X 分成可测集的有限分划的序列，使得分划 \(X_{n+1}\) 比 \(X_{n}\) 更细对所有 \(n \in N\) 成立，并且使得 \(X_{n}\) 的各部分的直径中的最大者当 \(n \to +\infty\) 时趋于 0 。
\end{enumerate}

设 \(E_{n} \subset E\) 是由满足如下条件的 \(f \in \mathrm{L}^{2}(\mathrm{X}, \mu)\) 组成的空间：f \(\mu\)-几乎处处在 \(X_{n}\) 的每个部分 Y 上为常数。我们令 \(E_{-1} = \{0\}\)。对每个 \(n \geqslant -1\)，用 \(p_{n}\) 表示 E 的以 \(E_{n}\) 为像的正交投影算子。

\begin{enumerate}
\def\labelenumi{\alph{enumi})}
\item
  序列 \((p_{n})_{n\in\mathbf{N}}\) 在赋有逐点收敛拓扑的 E 的自同态所成的空间 \(\mathcal{L}_{s}(\mathrm{E})\) 中收敛于恒等映射。（首先考虑 \(p_{n}(f)\) 当 \(f\in\mathcal{C}(\mathrm{X})\) 时的极限。）
\item
  对每个 \(n \in N\) ，自同态 \(q_{n} = p_{n} - p_{n-1}\) 是到 \(F_{n} = E_{n} \cap E_{n-1}^{\circ}\) 上的正交投影算子，并且诸空间 \(F_{n}\) 两两正交。
\item
  我们有 \(\sum_{n} q_{n} = 1_{\mathrm{E}}\) 在 \(\mathcal{L}_s(\mathrm{E})\) 中成立，并且空间 E 是诸空间 \(\mathrm{F}_n\) 的希尔伯特直和。
\end{enumerate}

\begin{enumerate}
\def\labelenumi{\arabic{enumi})}
\setcounter{enumi}{18}
\tightlist
\item
  设 \(\mathrm{X} \subset \mathbf{R}\) 是 \(\mathbf{R}\) 的紧子集，\(\mu\) 是 \(\mathrm{X}\) 上的正测度，而 \(u\) 是由 \(\mathrm{X}\) 的恒等函数给出的乘法自同态。
\end{enumerate}

我们记 \(\mathrm{E}=\mathrm{L}^{2}(\mathrm{X},\mu)\) ，并用 \(\mathcal{L}_{s}(\mathrm{E})\) 表示赋有逐点收敛拓扑的 E 的自同态所成的空间。

设 \(\varepsilon > 0\)。取 \(r > 0\) 使得 \(\mathrm{X} \subset ] - r, r[\) ，并取 \(\ell \geqslant 1\) 使得 \(r2^{-\ell} < \varepsilon\) 。对每个整数 \(n \in \mathbf{N}\) 以及每个整数 \(j\) ，其中 \(|j| \leqslant 2^{n + \ell}\) ，令

\[
\mathrm{X} _ {n, j} = \mathrm{X} \cap \Big ] \frac {r j}{2 ^ {n + \ell}}, \frac {r (j + 1)}{2 ^ {n + \ell}} \Big ]
\]

\begin{enumerate}
\def\labelenumi{\alph{enumi})}
\item
  对每个整数 \(n \in N\)，族 \(\mathcal{X}_{n} = (\mathrm{X}_{n,j})_{j}\) 是 X 的有限分划；诸集合 \(X_{n,j}\) 是可测的且直径 \(\leqslant r2^{-n-\ell}\) 。对每个 \(n \in N\)，分划 \(X_{n+1}\) 比分划 \(X_{n}\) 更细。因此我们可以对 \((\mathcal{X}_{n})\) 应用上一个习题的结果并沿用其记号，特别是空间 \(E_{n}\) 与 \(F_{n}\) 以及正交投影算子 \(p_{n}\) 与 \(q_{n}\) 。诸空间 \(E_{n}\) 是有限维的。
\item
  级数
\end{enumerate}

\[
v = \sum_ {n \in \mathbf {N}} q _ {n} u q _ {n}
\]

在空间 \(\mathcal{L}_{s}(\mathrm{E})\) 中收敛；存在 E 的标准正交基 B 使得 \(v \in \mathcal{D}_{\mathrm{B}}(\mathrm{E})\) 。

\begin{enumerate}
\def\labelenumi{\alph{enumi})}
\setcounter{enumi}{2}
\tightlist
\item
  自同态 \(w = v - u\) 是紧的。（考虑乘法自同态 \(u_n\) ，它由函数 \(g_n \in \mathrm{E}_n\) 给出，使得对每个整数 \(j\) （\(|j| \leqslant 2^{n + \ell}\) ）以及所有 \(x \in \mathrm{X}_{n,j}\) ，有 \(g_n(x) = rj2^{-n - \ell}\) ；然后验证 \(\| u - u_n \| \leqslant r2^{-n - \ell}\) 以及 \(u_n p_n = p_n u_n\) 。）
\end{enumerate}

\begin{enumerate}
\def\labelenumi{\arabic{enumi})}
\setcounter{enumi}{19}
\tightlist
\item
  设 E 是可数型复希尔伯特空间，\(u \in \mathcal{L}(\mathrm{E})\) 是埃尔米特自同态。设 \(\varepsilon\) 是严格正的实数。存在 E 的标准正交基 B、埃尔米特对角自同态 \(v \in \mathcal{D}_{\mathrm{B}}(\mathrm{E})\) 以及紧自同态 \(w \in \mathcal{L}^{\mathrm{c}}(\mathrm{E})\) ，使得 \(u = v + w\) 且 \(\|w\| < \varepsilon\) （首先考虑 u 有循环向量的情形，然后应用上一个习题。）
\end{enumerate}

21) 设 E 是可数型复希尔伯特空间，\(u \in \mathcal{L}(\mathrm{E})\) 是正规自同态。设 \(\varepsilon\) 是严格正的实数。存在 E 的标准正交基 B、对角自同态 \(v \in \mathcal{D}_{\mathrm{B}}(\mathrm{E})\) 以及紧自同态 \(w \in \mathcal{L}^{\mathrm{c}}(\mathrm{E})\) ，使得 \(u = v + w\) 且 \(\|w\| < \varepsilon\) 。

\begin{enumerate}
\def\labelenumi{\arabic{enumi})}
\setcounter{enumi}{21}
\tightlist
\item
  设 G 是局部紧交换拓扑群，\(\mu\) 是 G 上的哈尔测度。
\end{enumerate}

设 \(u \in \mathcal{L}(\mathrm{L}^2(\mathrm{G},\mu))\) 满足：\(u\) 与乘法算子 \(m_{\chi}\) 对每个特征标 \(\chi \in \widehat{\mathrm{G}}\) 交换。证明存在 \(g \in \mathrm{L}^{\infty}(\mathrm{G},\mu)\) 使得 \(u = m_g\) 。

\begin{enumerate}
\def\labelenumi{\arabic{enumi})}
\setcounter{enumi}{22}
\tightlist
\item
  设 \(u\) 是 E 的自同态，并设 A 是 \(\mathcal{L}(\mathrm{E})\) 的由 \(u\) 生成的闭含单位元子代数。
\end{enumerate}

若 v 是 E 的自同态且 \(u \circ v - v \circ u\) 属于 A，则 \(u \circ v - v \circ u\) 是拟幂零的。特别地，若 E 非零，我们有 \(u \circ v - v \circ u \neq 1_{E}\) 。（利用 I, p. 167 的习题 10；注意该结果对任何巴拿赫空间 E 都成立。）

\begin{enumerate}
\def\labelenumi{\arabic{enumi})}
\setcounter{enumi}{23}
\tightlist
\item
  设 \(\mathrm{A} \subset \mathcal{L}(\mathrm{E})\) 是 \(\mathcal{L}(\mathrm{E})\) 的含单位元的自伴子代数。我们用 \(\mathrm{A}''\) 表示 A 的双交换子。我们用 \(\mathcal{L}(\mathrm{E})_s\) 表示赋有逐点收敛拓扑的空间 \(\mathcal{L}(\mathrm{E})\) 。
\end{enumerate}

\begin{enumerate}
\def\labelenumi{\alph{enumi})}
\item
  空间 \(A''\) 是 \(\mathcal{L}(\mathrm{E})\) 的含单位元的自伴子代数，并且它在 \(\mathcal{L}(\mathrm{E})_s\) 中是闭的。
\item
  设 \(v \in A''\) 且 \(x \in E\)。向量 \(v(x)\) 属于诸 \(u(x)\) （\(u \in A\)）所成之集在 E 中的闭包。（设 F 是该闭包；注意 E 的以 F 为像的正交投影算子属于 \(A'\)。）
\item
  代数 \(A''\) 与 A 在 \(\mathcal{L}(\mathrm{E})_s\) 中的闭包重合。（对 E 的元素的每个有限族 \((x_i)_{i \in \mathrm{I}}\) ，把上一问题应用于空间 \(\mathrm{E}^{\mathrm{I}}\) 、像 \(\widetilde{\mathrm{A}}\) （即 A 在含单位元对合态射 \(f: \mathcal{L}(\mathrm{E}) \to \mathcal{L}(\mathrm{E}^{\mathrm{I}})\) （其中 \(f(u)(y_i) = (u(y_i))_{i \in \mathrm{I}}\) ）下的像），以及元素 \((x_i)\) （它属于 \(\mathrm{E}^{\mathrm{I}}\) ）。）
\end{enumerate}

\(A''\) 与 A 在 \(\mathcal{L}(\mathrm{E})_s\) 中闭包的相等就是冯·诺伊曼双交换子定理。

\begin{enumerate}
\def\labelenumi{\arabic{enumi})}
\setcounter{enumi}{24}
\tightlist
\item
  设 U 是 R 的开子集。普遍可测函数 \(f: U \to R\) 称为勒夫纳函数，如果对每个复希尔伯特空间 E 以及 E 的满足 \(\operatorname{Sp}(u) \subset \operatorname{U}\) 与 \(\operatorname{Sp}(v) \subset \operatorname{U}\) 的埃尔米特自同态 u 与 v，条件 \(u \leqslant v\) 在 \(\mathcal{L}(\operatorname{E})\) 中蕴含 \(f(u) \leqslant f(v)\) 。我们用 \(\operatorname{Loe}(\operatorname{U})\) 表示 U 上勒夫纳函数所成的集。
\end{enumerate}

\begin{enumerate}
\def\labelenumi{\alph{enumi})}
\item
  设 \(f \in \text{Loe}(U)\)。函数 f 在 U 上递增。常值函数 1 与 U 的恒等函数都属于 \(\text{Loe}(U)\) 。
\item
  设 \(f \in \text{Loe}(\mathbf{R}_+^*)\)。设 E 是复希尔伯特空间，u 是 E 的正自同态。对 E 的每个正交投影算子 p 以及每个 \(\varepsilon > 0\) ，我们有 \(f(u)_p \leqslant f((1 + \varepsilon)u_p)\) ，其中用 \(v_p\) 表示 \(\text{Im}(p)\) 的由 \(x \mapsto p(v(x))\) 定义的自同态（对所有 \(v \in \mathcal{L}(\text{E})\) ）。（把 E 等同于希尔伯特直和 \(\text{Im}(p) \oplus \text{Im}(1_{\text{E}} - p)\) ，并将 u 表示为矩阵 \(\begin{pmatrix} a & b \\ c & d \end{pmatrix}\)。注意 \(u \leqslant \begin{pmatrix} (1 + \varepsilon)a & 0 \\ 0 & (1 + \varepsilon^{-1})d \end{pmatrix}\) 对所有 \(\varepsilon > 0\) 成立。）
\item
  设 \(f \in \text{Loe}(\mathbf{R}_{+}^{*})\)。函数 f 在 \(R_{+}^{*}\) 上是凹的（设 v 是 \(C^{2}\) 的在典范基下为对角的自同态，并设 \(g \in \text{SO}(\mathbf{R}^{2})\) 视为 \(C^{2}\) 的自同态。把上一问题应用于以 \(Ce_{1}\) 为像的正交投影算子且取 \(u = g^{*}vg\) 。）
\item
  集合 \(\operatorname{Loe}(\mathbf{R})\) 是形如 \(f(x)=ax+b\) （其中 \((a,b)\in\mathbf{R}_{+}\times\mathbf{R}\) ）的函数所成的集。（把上一问题既应用于 f，也应用于函数 \(x\mapsto-f(-x)\) 。）
\item
  设 \(t \in \mathbf{R}_{+}^{*}\) 。函数 \(f: x \mapsto -(t + x)^{-1}\) 属于 \(\operatorname{Loe}(\mathbf{R}_{+}^{*})\)。
\end{enumerate}

\(f)\) 对任何有界正测度 \(\nu\) （在 \(\mathbf{R}_{+}^{*}\) 上且使函数 \(x\mapsto x^{-1}\) 为 \(\nu\) -可积），函数 \(f\colon\mathbf{R}_{+}^{*}\to\mathbf{R}\) 由

\[
f (x) = \int_ {\mathbf {R} _ {+} ^ {*}} \frac {x}{x + t} d \nu (t)
\]

所给出的函数属于 \(\text{Loe}(\mathbf{R}_{+}^{*})\)；它在 \(R_{+}^{*}\) 上可微，并且还满足下列条件：

\begin{enumerate}
\def\labelenumi{(\roman{enumi})}
\item
  \(\lim_{t\to0}f(t)=0;\)
\item
  \(\lim_{t\to0}f'(t)\) 在 R 中存在；
\item
  f 在 \(+\infty\) 处有有限极限。
\end{enumerate}

\begin{enumerate}
\def\labelenumi{\arabic{enumi})}
\setcounter{enumi}{25}
\tightlist
\item
  设 E 是复希尔伯特空间，a 是 E 的正自同态。我们记 \(\|x\|_{a} = \langle x | a(x) \rangle\) ；它是 E 上的半范数。我们同样用 \(\|u\|_{a}\) 表示 E 赋有该半范数时其自同态 u 的半范数。
\end{enumerate}

我们称函数 \(f \colon \mathbf{R}_+^* \to \mathbf{R}_+^*\) 为插值函数，如果对每个复希尔伯特空间 E、E 的每个正自同态 \(a\) 以及每个 \(u \in \mathcal{L}(\mathrm{E})\) ，条件 \(\| u \| \leqslant 1\) 与 \(\| u \|_a \leqslant 1\) 蕴含 \(\| f(u) \|_a \leqslant 1\) 。

\begin{enumerate}
\def\labelenumi{\alph{enumi})}
\item
  函数 \(f \colon \mathbf{R}_+^* \to \mathbf{R}_+^*\) 是插值函数，当且仅当对每个复希尔伯特空间 E、E 的每个正自同态 \(a\) 以及每个 \(u \in \mathcal{L}(\mathrm{E})\) ，条件 \(u^*u \leqslant 1_{\mathrm{E}}\) 与 \(u^*au \leqslant a\) 蕴含 \(u^*f(a)u \leqslant f(a)\) 。
\item
  \(R_{+}^{*}\) 上的任何勒夫纳函数都是插值函数。（给定 E 以及 \(u \in \mathcal{L}(E)\) （范数 \(\leqslant 1\) ），令 \(w = \sqrt{1_{E} - u^{*}u}\) ，\(w' = \sqrt{1_{E}-uu^{*}}\) ；考虑 \(F = E \oplus E\) ，并注意由 \(v(x,y) = (u(x)+w'(y), -w(x)+u(y))\) 定义的 F 的自同态 v 是等距的；然后考虑 \(a' \colon (x,y) \mapsto (a(x),0)\) ，并把上一个习题的 b) 应用于以 \(E \oplus \{0\}\) 为像的正交投影算子。）
\item
  每个插值函数都是 \(\mathbf{R}_{+}^{*}\) 上的勒夫纳函数。（设 \(0 \leqslant u \leqslant v\) 是复希尔伯特空间 E 的自同态且 \(v\) 可逆；令 \(\mathrm{F} = \mathrm{E} \oplus \mathrm{E}\) ，考虑 F 的自同态 \(a = u \oplus v\) 以及自同态 \(w\) ，它由矩阵 \(\left( \begin{array}{cc}0 & 1_{\mathrm{E}}\\ v^{-1/2}u^{1/2} & 0 \end{array} \right)\) 表示；证明 \(w^{*}w \leqslant 1_{\mathrm{F}}\) 且 \(w^{*}aw \leqslant a\) 。）
\item
  设 f 是插值函数。设 E 是复希尔伯特空间；我们记 \(\mathcal{R}(u)=\frac{1}{2}(u+u^{*})\) 对所有 \(u\in\mathcal{L}(\mathrm{E})\) 。对 \(\mathcal{L}(\mathrm{E})\) 中的 u 与 v，证明：\(e^{-tv^{*}}\circ u\circ e^{-tv}\leqslant u\) 对所有 \(t\in R_{+}\) 成立当且仅当 \(\mathcal{R}(u\circ v)\geqslant0\) 。（对给定的 \(x\in E\)，计算函数 \(t\mapsto\langle e^{-tv}x|(ue^{-tv})x\rangle\) 的导数。）
\end{enumerate}

由此推出：对 E 的所有自同态 \(u\) 与 \(v\) ，当条件

\[
u \geqslant 0, \quad \mathcal {R} (v) \geqslant 0, \quad \mathcal {R} (u v) \geqslant 0
\]

得到满足时，我们有 \(\mathcal{R}(f(u)v) \geqslant 0\) 。

27) 我们打算证明习题 25 问题 \(b\) 中论断的逆命题。为此，将利用下面 V, p. 503 习题 23 的结果与记号。特别地，G 表示群 \(\mathbf{R}_{+}^{*}\) ，它赋有哈尔测度 \(\mu = x^{-1}dx\) ，并与群 \(i\mathbf{R}\) 通过映射 \((x, it) \mapsto x^{it}\) 对偶。我们用 X 表示 G 的恒等映射。

我们分别用 co、si 与 ta 表示 iR 上的函数 \(it \mapsto \cos(i\pi t)\)、\(it \mapsto \sin(i\pi t)\) 与 \(it \mapsto \tan(i\pi t)\) 。

设 \(f \in \text{Loe}(\mathbf{G})\) 满足习题 25 问题 \(f\) 的条件 (i)、(ii) 与 (iii)。

\begin{enumerate}
\def\labelenumi{\alph{enumi})}
\item
  设 \(\alpha\) 是由 \(x \mapsto x / (1 + x)\) 给出的 \(G\) 上的函数。我们有 \(\mathrm{I}_{\alpha} = ] - 1,0[\) ，并且 \(\widehat{\alpha}(s) = -\pi / \sin(\pi s)\) 对 \(s \in \mathrm{B}_f\) 成立。
\item
  函数 ta 在 \(i\mathbf{R}\) 上有界。我们令 \(v = F \circ m_{ta} \circ \overline{F}\) ；它是 \(L^{2}(G)\) 的自同态，满足 \(\|v\| \leqslant 1\) 且 \(\mathcal{R}(v) = 0\) 。
\item
  设 \(\varepsilon > 0\) 是实数。我们用 \(u_{\varepsilon}\) 表示由 \(h_{\varepsilon} = \mathrm{X}/(1 + \varepsilon\mathrm{X})\) 给出的 \(\mathrm{L}^2(\mathrm{G})\) 上的乘法自同态。我们有 \(u_{\varepsilon} \geqslant 0\) ，并且对每个 \(\varphi \in \mathrm{L}^2(\mathrm{G})\) 以及每个有界连续函数 \(f\) （在 \(\mathrm{G}\) 上），有
\end{enumerate}

\[
\langle \varphi | \mathcal {R} (f (u _ {\varepsilon}) v) \varphi \rangle = \frac {1}{2} \langle f \circ h _ {\varepsilon} | \mathrm{Q} (\varphi) \rangle
\]

\(\mathrm{Q}(\varphi) = \overline{\varphi} v(\varphi) + \overline{v(\varphi)}\varphi .\)

\begin{enumerate}
\def\labelenumi{\alph{enumi})}
\setcounter{enumi}{3}
\item
  设 F 是 \(\mathrm{L}^{2}(i\mathbf{R})\) 的由诸函数 \(t \mapsto e^{\alpha(it - s_{0})^{2}}\) （其中 \(\alpha \in R_{+}^{*}\) ，\(s_{0} \in C\) ）生成的向量子空间；它在 \(\mathrm{L}^{2}(i\mathbf{R})\) 中稠密。
\item
  设 \(\psi \in \mathrm{F}\) 。存在 \(\varphi \in \mathrm{L}^2 (\mathbf{G})\) 使得 \(\widehat{\varphi} = \mathrm{co}\cdot \psi\)。令 \(g = \mathrm{Q}(\varphi)\)。我们有 \(g\in \mathrm{L}^{1}(\mathrm{G})\) ，并且存在 \(\mathbf{C}\) 上的全纯函数，它到 \(i\mathbf{R}\) 的限制与 \(\widehat{g}\) 一致。此外
\end{enumerate}

\[
\widehat {g} = \frac {1}{2 i \pi} (\psi^ {*} * \psi) \mathrm{si}.
\]

\(f)\) 对每个 \(\varepsilon > 0\) ，我们有 \(\mathcal{R}(u_{\varepsilon} \circ v) \geqslant 0\) 。（计算 \(\langle \varphi | \mathcal{R}(u_{\varepsilon} u) \varphi \rangle\) ，当 \(\varphi\) 是上一问题中那样的函数时；并像 V, p. 503 习题 23 的问题 \(d\) 中那样使用 G 上的卷积。）

\(g)\) 设 \(f \in \mathrm{Loe}(\mathbf{R}_{+}^{*})\)。设 \(g\) 是与 \(\psi \in F\) 如问题 \(e)\) 中那样相伴的函数。我们则有

\[
\int_ {\mathrm{G}} f g d \mu \geqslant 0
\]

以及

\[
\int_ {\mathbf {R}} \widehat {f} (- \sigma - i t) \sin (\pi (\sigma + i t)) (\psi^ {*} * \psi) (\sigma + i t) d t \geqslant 0
\]

对所有 \(\sigma\in]0,1[\) 成立。

\begin{enumerate}
\def\labelenumi{\alph{enumi})}
\setcounter{enumi}{7}
\tightlist
\item
  设 \(\sigma \in ]0,1[\) 。函数 \(q\) 由 \(q(s) = \widehat{f}(-s)\sin(\pi s)\) （对 \(s\) 满足 \(\mathcal{R}(s) = \sigma\) ）定义，它满足
\end{enumerate}

\[
\sum_ {i = 0} ^ {n} \sum_ {j = 0} ^ {n} \bar {t} _ {i} t _ {j} q \left(\frac {1}{2} (\bar {z} _ {i} + z _ {j})\right) \geqslant 0
\]

对所有 \(n \in N\) 以及所有复数族 \((t_{i})_{0 \leqslant i \leqslant n}\) 与 \((z_{i})_{0 \leqslant i \leqslant n}\) （它们满足对每个 i 有 \(\mathcal{R}(z_{i}) = \sigma\) ）成立。

\begin{enumerate}
\def\labelenumi{\roman{enumi})}
\tightlist
\item
  存在 G 上的正测度 \(\nu_0\) ，使得 \(\widehat{\nu}_0(it) = \widehat{f}(-it)\sin(i\pi t)/\pi\) 对所有 \(t\in\mathbf{R}\) 成立；我们有 \(\mathrm{X}^{\sigma}\in\mathrm{L}^{1}(\mathrm{G},\nu_{0})\) ，若 \(0<\sigma<1\) 。
\end{enumerate}

\begin{enumerate}
\def\labelenumi{\alph{enumi})}
\setcounter{enumi}{9}
\tightlist
\item
  设 \(\nu = \check{\nu}_0\)。测度 \(\nu\) 是有界的。（先验证 \(]-1,0[ \subset I_{\nu}\) ，再验证 \(\widehat{\nu}(-t)\) 当 \(t \to 0\) 时有界，以此证明 \(\nu([1,+\infty[)\) 有限。）对每个 \(x\) ，我们有
\end{enumerate}

\[
f (x) = \int_ {\mathrm{G}} \frac {x}{x + t} d \nu (t).
\]

\begin{enumerate}
\def\labelenumi{\arabic{enumi})}
\setcounter{enumi}{27}
\tightlist
\item
  \begin{enumerate}
  \def\labelenumii{\alph{enumii})}
  \tightlist
  \item
    设 \(D \subset C\) 是开单位圆盘。设 f 是 \(\overline{D}\) 的一个邻域中的全纯函数的芽。对每个 \(z \in D\) ，我们有
  \end{enumerate}
\end{enumerate}

\[
f (z) = i \mathcal {I} (f (0)) + \frac {1}{2 \pi} \int_ {0} ^ {2 \pi} \mathcal {R} (f (e ^ {i \theta})) k (e ^ {i \theta}, z) d \theta
\]

其中 \(k(z, w) = (z + w) / (z - w)\) 对 \((z, w) \in \mathrm{D} \times \mathrm{D}\) 。（注意

\[
k (e ^ {i \theta}, z) = 1 + 2 \sum_ {n = 1} ^ {+ \infty} z ^ {n} e ^ {- i n \theta}
\]

并使用 f 在 0 附近的泰勒展开。）

\begin{enumerate}
\def\labelenumi{\alph{enumi})}
\setcounter{enumi}{1}
\tightlist
\item
  设 \(f\) 是 D 上的全纯函数，使得 \(\mathcal{R}(f(z)) > 0\) 对所有 \(z \in \mathrm{D}\) 成立。测度 \(\mu_r = \frac{1}{2\pi} \mathcal{R}(f(re^{i\theta})) d\theta\) （\([0, 2\pi]\) 上的）当 \(r \to 1\) 时淡收敛于一个有界正测度 \(\nu\) ，并且我们有
\end{enumerate}

\[
f (z) = i \mathcal {I} (f (0)) + \frac {1}{2 \pi} \int_ {0} ^ {2 \pi} \mathcal {R} (f (e ^ {i \theta})) k (e ^ {i \theta}, z) d \nu (\theta)
\]

对所有 \(z \in D\) 成立。（化归到 \(f(0) = 1\) 的情形，并注意这时诸测度 \(\mu_{r}\) 是总质量为 1 的正测度。）

\begin{enumerate}
\def\labelenumi{\alph{enumi})}
\setcounter{enumi}{2}
\tightlist
\item
  设 H 是由这样的 \(z \in C\) 组成的 C 的开子集：\(\mathcal{I}(z) > 0\)。设 \(f: H \to C\) 是全纯函数。f 的像含于 H 当且仅当存在 R 上的有界正测度 \(\nu\) 以及实数 \(a \geqslant 0\) 使得
\end{enumerate}

\[
f (z) = \mathcal {R} (f (i)) + a z + \int_ {\mathbf {R}} \frac {1 + x z}{x - z} d \nu (x)
\]

对所有 \(z \in \mathbf{H}\) 成立。

\begin{enumerate}
\def\labelenumi{\alph{enumi})}
\setcounter{enumi}{3}
\tightlist
\item
  设 I 是 \(\mathbf{R}\) 中的非空开区间，并设 \(f\colon \mathrm{I}\to \mathbf{R}\) 。记 \(\mathrm{J} = \mathbf{R} - \mathrm{I}\) 。下列条件等价（“勒夫纳定理”）：
\end{enumerate}

\begin{enumerate}
\def\labelenumi{(\roman{enumi})}
\item
  函数 \(f\) 属于 \(\operatorname{Loe}(\mathbf{I})\) ；
\item
  存在解析函数 \(\widetilde{f}\) ，它定义在 \((\mathbf{C} - \mathbf{R}) \cup I\) 上、到 I 的限制与 \(f\) 一致，并且满足 \(\mathcal{I}(f(z)) > 0\) 若 \(\mathcal{I}(z) > 0\) ；
\item
  存在 J 上的有界测度 \(\nu\) 以及偶 \((a,b)\in\mathbf{R}_{+}\times\mathbf{R}\) 使得
\end{enumerate}

\[
f (x) = a x + b + \int_ {\mathrm{J}} \frac {1 + x y}{y - x} d \nu (y)
\]

对所有 \(x \in I\) 成立。

（首先直接验证 (iii) 蕴含 (i) 与 (ii)；为证明 (i) 蕴含 (iii)，单独处理 I = R 的情形，并把其他情形化归为 I = R \(_{+}^{*}\) 的情形，对此应用上一个习题；最后，为证明 (ii) 蕴含 (iii)，应用 c) 部分并验证由此给出的 R 上的测度 \(\nu\) 的支撑含于 J。）

\exercisesection{}

\begin{enumerate}
\def\labelenumi{\arabic{enumi})}
\item
  函数 \(\varphi_0\) 在 \(\mathbf{R}\) 上由 \(\varphi_0(x) = 0\) （若 \(x \leqslant 0\) ）与 \(\varphi_0(x) = e^{-1/x}\) （若 \(x > 0\) ）定义，它是 \(\mathrm{C}^\infty\) 类的，并且满足 \(\varphi_0^{(k)}(0) = 0\) 对每个整数 \(k \in \mathbf{N}\) 。
\item
  设 \(x_0 \in \mathbf{R}^n\) ，\(r > 0\) 且 \(t > 1\) 。存在 \(\varphi \in \mathcal{D}(\mathbf{R}^n)\) 使得
\end{enumerate}

\begin{enumerate}
\def\labelenumi{(\roman{enumi})}
\item
  我们有 \(0 \leqslant \varphi \leqslant 1\) ；
\item
  \(\varphi(x)=1\) 对 \(\|x-x_{0}\|\leqslant r\) 成立，且 \(\varphi(x)=0\) 对 \(\|x-x_{0}\|\geqslant tr\) 成立；
\item
  对每个 \(\alpha \in \mathbf{N}^n\) ，我们有
\end{enumerate}

\[
\sup _ {x \in \mathbf {R} ^ {n}} | \partial^ {\alpha} \varphi (x) | \leqslant \alpha ! ((t - 1) r) ^ {| \alpha |}.
\]

\begin{enumerate}
\def\labelenumi{\arabic{enumi})}
\setcounter{enumi}{2}
\tightlist
\item
  设 \(n \in \mathbf{N}\) 。我们用 \(\mu_n\) 表示 \(\mathbf{R}^n\) 上的勒贝格测度。
\end{enumerate}

\begin{enumerate}
\def\labelenumi{\alph{enumi})}
\tightlist
\item
  设 \(\mathcal{H}_1\) 是映射 \(\partial_1\colon \mathcal{D}(\mathbf{R}^n)\to \mathcal{D}(\mathbf{R}^n)\) 的像；它是满足如下条件的 \(\varphi \in \mathcal{D}(\mathbf{R}^n)\) 所成的空间：
\end{enumerate}

\[
\int_ {\mathbf {R}} \varphi (x, y) d \mu_ {1} (x) = 0
\]

对所有 \(y \in \mathbf{R}^{n-1}\) 成立。

\begin{enumerate}
\def\labelenumi{\alph{enumi})}
\setcounter{enumi}{1}
\tightlist
\item
  设 \(\varphi_0\in \mathcal{D}(\mathbf{R})\) 使其在 \(\mathbf{R}\) 上的积分等于 1。对每个 \(\varphi \in \mathcal{D}(\mathbf{R}^n)\) ，存在 \(\varphi_{1}\in \mathcal{H}_{1}\) 使得
\end{enumerate}

\[
\varphi (x, y) = \varphi_ {1} (x, y) + \varphi_ {0} (x) \int_ {\mathbf {R}} \varphi (t, y) d \mu_ {1} (t)
\]

对所有 \((x,y)\in \mathbf{R}\times \mathbf{R}^{n - 1}\) 成立。

\begin{enumerate}
\def\labelenumi{\alph{enumi})}
\setcounter{enumi}{2}
\tightlist
\item
  对每个广义函数 \(f \in \mathcal{D}'(\mathbf{R}^{n})\) ，存在 \(g \in \mathcal{D}'(\mathbf{R}^{n})\) 使得 \(\partial_{1}g = f\) 。若 \(g_{1}\) 与 \(g_{2}\) 属于 \(\mathcal{D}'(\mathbf{R}^{n})\) 且满足 \(\partial_{1}g_{1} = \partial_{1}g_{2} = f\) ，则存在广义函数 \(h \in \mathcal{D}'(\mathbf{R}^{n-1})\) 使得
\end{enumerate}

\[
\langle g _ {1} - g _ {2}, \varphi \rangle = \langle h, \psi \rangle
\]

对所有 \(\varphi\in\mathcal{D}(\mathbf{R}^{n})\) 成立，其中 \(\psi\in\mathcal{D}(\mathbf{R}^{n-1})\) 由下式定义：

\[
\psi (y) = \int_ {\mathbf {R}} \varphi (t, y) d \mu_ {1} (t).
\]

\(d)\) 设 U 是 \(\mathbf{R}^{n}\) 的连通开子集。设 \(f \in \mathcal{D}'(\mathrm{U})\) 是满足 \(\partial_{i}f = 0\) （对 \(1 \leqslant i \leqslant n\) ）的广义函数。广义函数 \(f\) 是与一个常值函数相伴的广义函数。

\begin{enumerate}
\def\labelenumi{\arabic{enumi})}
\setcounter{enumi}{3}
\tightlist
\item
  我们用 \(\mu\) 表示 \(\mathbf{R}\) 上的勒贝格测度。
\end{enumerate}

\begin{enumerate}
\def\labelenumi{\alph{enumi})}
\tightlist
\item
  设 \(\nu\) 是 \(\mathbf{R}\) 上的测度，\(a \in \mathbf{R}\) 且 \(\nu(\{a\}) = 0\) 。对 \(x \in \mathbf{R}\) 我们令
\end{enumerate}

\[
f (x) = \left\{ \begin{array}{l l} \nu ([ a, x ]) & \text { if } x \geqslant a, \\ \nu ([ x, a ]) & \text { if } x \leqslant a. \end{array} \right.
\]

函数 f 是连续的，并且与之相伴的广义函数满足 \(f' = \nu\) 。

\begin{enumerate}
\def\labelenumi{\alph{enumi})}
\setcounter{enumi}{1}
\item
  设 f 是 R 上的广义函数，使得 \(f'\) 是一个测度 \(\nu\) 。广义函数 f 与 R 上按上一问题方式定义的连续函数相伴。
\item
  设 \(f\) 是 \(\mathbf{R}\) 上的广义函数，使得 \(f^{(k)}\) 对每个整数 \(k \geqslant 0\) 都是测度。那么 \(f\) 是与 \(\mathcal{C}^{\infty}(\mathbf{R})\) 中的函数相伴的广义函数。
\end{enumerate}

\begin{enumerate}
\def\labelenumi{\arabic{enumi})}
\setcounter{enumi}{4}
\item
  与点测度 \(\varepsilon_{n}\) （\(n\in\mathbf{Z}\) ）相伴的广义函数族在 \(\mathcal{S}^{\prime}(\mathbf{R})\) 中可和。其和 \(\omega\) 满足 \(\mathcal{F}(\omega)=\omega\) 。（该公式等价于关于赋有勒贝格测度的局部紧群 \(\mathrm{G}=\mathbf{R}^{n}\) 与子群 \(\mathrm{H}=\mathbf{Z}^{n}\) 的泊松公式，参见 II, § 1, 第 7 目。）
\item
  设 A 是 C 上的交换含单位元代数，满足下列条件：
\end{enumerate}

\begin{enumerate}
\def\labelenumi{(\roman{enumi})}
\item
  代数 \(\mathcal{C}(\mathbf{R})\) 是 A 的含单位元子代数；
\item
  存在导子 \(d\colon\mathrm{A}\to\mathrm{A}\) ，它到 \(\mathcal{C}^{1}(\mathbf{R})\) 的限制与函数的通常导子一致。
\end{enumerate}

\begin{enumerate}
\def\labelenumi{\alph{enumi})}
\item
  R 的恒等映射 X 在 A 中可逆。（证明映射 \(g: \mathbf{R} \to \mathbf{R}\) （它满足 \(g(0) = 0\) ，并且 \(g(x) = x(\log(|x|) - 1)\) 对 \(x \neq 0\) ）是连续的且满足 \(gX = 1\) 。）
\item
  我们有 \(d \circ d(|\mathrm{X}|) = 0\) 。
\item
  不存在双线性映射 \(m\colon\mathcal{D}^{\prime}(\mathbf{R})\times\mathcal{D}^{\prime}(\mathbf{R})\to\mathcal{D}^{\prime}(\mathbf{R})\) 使得 \(m(f,g)=fg\) 对每个偶 \((f,g)\in\mathcal{D}(\mathbf{R})^{2}\) 成立，并且使得
\end{enumerate}

\[
m (f, g) ^ {\prime} = m \left(f ^ {\prime}, g\right) + m \left(f, g ^ {\prime}\right)
\]

对所有 \((f,g)\in\mathcal{D}^{\prime}(\mathbf{R})^{2}\) 成立（“广义函数乘法的不可能性”）。

\begin{enumerate}
\def\labelenumi{\alph{enumi})}
\setcounter{enumi}{3}
\tightlist
\item
  给出 R 上试验函数的序列 \((f_{n})\) 与 \((g_{n})\) 的例子，使得 \(f_{n} \to \varepsilon_{0}\) 且 \(g_{n} \to \varepsilon_{0}\) 在 \(\mathcal{D}'(\mathbf{R})\) 中，并且 \(f_{n}^{2}\) 与 \(g_{n}^{2}\) 在 \(\mathcal{D}'(\mathbf{R})\) 中收敛于不同的极限。
\end{enumerate}

\begin{enumerate}
\def\labelenumi{\arabic{enumi})}
\setcounter{enumi}{6}
\tightlist
\item
  设 \(n \in N\) 。我们用 \(\widetilde{\gamma}\) 表示 \(R^{n}\) 在从 \(R^{n}\) 到 C 的函数所成空间上的线性表示，它由下式定义：
\end{enumerate}

\[
\widetilde {\gamma} (h) f (x) = f (x + h)
\]

对所有 \(h \in \mathbf{R}^n\) 以及每个函数 \(f: \mathbf{R}^n \to \mathbf{C}\) 成立。

\begin{enumerate}
\def\labelenumi{\alph{enumi})}
\item
  对每个 \(h \in R^{n}\) ，映射 \(\widetilde{\gamma}(h)\) 通过向子空间的过渡定义了局部凸空间 \(\mathcal{D}(\mathbf{R}^{n})\) 的一个自同构，它仍记作 \(\widetilde{\gamma}(h)\) 。
\item
  我们用 \(\pmb{\gamma}\) 表示 \(\mathcal{D}'(\mathbf{R}^n)\) 的与 \(\widetilde{\pmb{\gamma}}(-h)\) 转置相应的自同构。若 \(f \in \mathcal{D}(\mathbf{R}^n)\) ，则 \(\pmb{\gamma}(h)f = \widetilde{\pmb{\gamma}}(h)f\) 对所有 \(h \in \mathbf{R}^n\) 成立。
\item
  映射 \(h \mapsto \gamma(h)\) 是 \(\mathbf{R}^{n}\) 在 \(\mathcal{D}'(\mathbf{R}^{n})\) 中的线性表示。
\end{enumerate}

\(d)\) 对每个 \(h \in \mathbf{R}^n\) 以及每个 \(\alpha \in \mathbf{N}^n\) ，我们有 \(\partial^\alpha \circ \boldsymbol{\gamma}(h) = \boldsymbol{\gamma}(h) \circ \partial^\alpha\) 。

\begin{enumerate}
\def\labelenumi{\alph{enumi})}
\setcounter{enumi}{4}
\tightlist
\item
  设 i 是满足 \(1 \leqslant i \leqslant n\) 的整数，\(e_{i}\) 是第 \(i\) 个向量，属于 \(R^{n}\) 的典范基。对每个广义函数 \(f \in \mathcal{D}'(\mathbf{R}^{n})\) ，我们有
\end{enumerate}

\[
\lim_{\substack{h\to 0\\ h\neq 0}}\frac{1}{h} (\boldsymbol {\gamma}(he_{i})f - f) = \partial_{i}f
\]

在赋有弱拓扑的 \(\mathcal{D}'(\mathbf{R}^n)\) 中成立。

\begin{enumerate}
\def\labelenumi{\arabic{enumi})}
\setcounter{enumi}{7}
\item
  设 \(n \in \mathbf{N}\) ，U 是 \(R^{n}\) 的开子集。广义函数 \(f \in \mathcal{D}'(U)\) 称为正的，如果 \(\langle f, \varphi \rangle \geqslant 0\) 对每个 \(\varphi \in \mathcal{D}(U)\) （它满足 \(\varphi \geqslant 0\) ）成立。证明 f 是正的当且仅当 f 是与 U 上一个正测度相伴的广义函数。
\item
  设 \(n \in \mathbf{N}\) ，\(U\) 是 \(\mathbf{R}^n\) 的开子集。设 \(k \in \mathbf{N}\) 且 \(p \in ]1, +\infty[\) 。空间 \(W^{k,p}(U)\) 是自反的。
\item
  设 \(n \in \mathbf{N}\) ，\(U\) 是 \(\mathbf{R}^n\) 的开子集。设 \(k \in \mathbf{N}\) 且 \(p \in [1, +\infty[\)。设 \(V \subset U\) 是相对紧的开子集。
\end{enumerate}

设 \(\varphi \in \mathcal{D}(\mathbf{R}^n)\) 是正函数，其支撑含于 \(\mathbf{R}^n\) 的单位球且积分为 1。对每个 \(\varepsilon > 0\) 以及 \(x \in \mathbf{R}^n\) ，我们令 \(\varphi_{\varepsilon}(x) = \varepsilon^{-n}\varphi(x/\varepsilon)\)。设 \(f \in \mathrm{W}^{k,p}(\mathrm{U})\) 。

\begin{enumerate}
\def\labelenumi{\alph{enumi})}
\item
  \(f_{\mathrm{V}}\) （即 \(f\) 到 V 的限制）属于 \(\mathrm{W}^{k,p}(\mathrm{V})\) 。
\item
  设 \(\widetilde{f}\) 是 \(\mathbf{R}^n\) 上在 U 上与 \(f\) 一致且在 U 外为零的函数。当 \(\varepsilon \to 0\) 时，\(\widetilde{f} * \varphi_{\varepsilon}\) 到 V 的限制属于 \(\mathrm{W}^{k,p}(\mathrm{V})\) 且收敛于 \(f_{\mathrm{V}}\) 在 \(\mathrm{W}^{k,p}(\mathrm{V})\) 中。
\end{enumerate}

\begin{enumerate}
\def\labelenumi{\arabic{enumi})}
\setcounter{enumi}{10}
\tightlist
\item
  设 \(n \in \mathbf{N}\) ，U 是 \(\mathbf{R}^n\) 的开子集。设 \(k \in \mathbf{N}\) 且 \(p \in [1, +\infty[\) 。我们用 \(\mathrm{H}^{k,p}(\mathrm{U})\) 表示 \(\mathcal{C}^{\infty}(\mathrm{U})\) 在 \(\mathrm{W}^{k,p}(\mathrm{U})\) 中的闭包。
\end{enumerate}

对 \(x \in \mathrm{U}\) ，用 \(\delta(x)\) 表示从 \(x\) 到 \(\mathrm{U}\) 在 \(\mathbf{R}^n\) 中的边界的距离；令 \(\mathrm{U}_{-1} = \mathrm{U}_0 = \varnothing\) 并且

\[
\mathrm{U} _ {k} = \{x \in \mathrm{U} | \| x \| <   k, \delta (x) > k ^ {- 1} \}
\]

对每个整数 \(k \geqslant 1\) 。令 \(\mathrm{V}_k = \mathrm{U}_{k+1} \cap (\mathbf{R}^n - \overline{\mathrm{U}}_{k-1})\) 对 \(k \geqslant 1\)。诸开集 \(\mathrm{V}_k\) 覆盖 U；设 \((\psi_k)_{k \geqslant 1}\) 是从属于覆盖 \((\mathrm{V}_k)_{k \geqslant 1}\) 且由 \(\mathcal{D}(\mathrm{U})\) 中的函数组成的单位分解。

\begin{enumerate}
\def\labelenumi{\alph{enumi})}
\tightlist
\item
  对 \(k \geqslant 1\) ，函数
\end{enumerate}

\[
\varphi_{k} = \sum_{\substack{j\geqslant 1\\ \operatorname{Supp}(\psi_{j})\subset \mathrm{V}_{k}}}\psi_{j}
\]

属于 \(\mathcal{D}(\mathrm{U})\) 且支撑含于 \(\mathrm{V}_k\) ；我们有

\[
\sum_ {k \geqslant 1} \varphi_ {k} = 1
\]

在 U 上成立。

\begin{enumerate}
\def\labelenumi{\alph{enumi})}
\setcounter{enumi}{1}
\tightlist
\item
  设 \(f \in \mathrm{W}^{k,p}(\mathrm{U})\) 且 \(\varepsilon > 0\)。设 \(\varphi\) 以及 \(\varphi_{\varepsilon}\) （对 \(\varepsilon > 0\) ）是如习题 10 中那样的函数。对每个 \(k \geqslant 1\)，存在 \(\varepsilon_{k} > 0\) 使得
\end{enumerate}

\[
\left\| \varphi_ {\varepsilon_ {k}} * (\psi_ {k} f) - \psi_ {k} f \right\| _ {k, p} <   \frac {\varepsilon}{2 ^ {k}}.
\]

\begin{enumerate}
\def\labelenumi{\alph{enumi})}
\setcounter{enumi}{2}
\tightlist
\item
  函数
\end{enumerate}

\[
f _ {\varepsilon} = \sum_ {k \geqslant 1} \varphi_ {\varepsilon_ {k}} * (\psi_ {k} f)
\]

属于 \(\mathcal{C}^{\infty}(\mathrm{U})\) 且 \(\| f - f_{\varepsilon}\|_{k,p} < \varepsilon\) 。

\(d)\). 空间 \(\mathrm{H}^{k,p}(\mathrm{U})\) 等于 \(\mathrm{W}^{k,p}(\mathrm{U})\) 。

\begin{enumerate}
\def\labelenumi{\alph{enumi})}
\setcounter{enumi}{4}
\item
  空间 \(\mathrm{H}^{k,p}(\mathrm{U})\) 可等同于 \(\mathcal{C}^{k}(\mathrm{U})\) 关于范数 \(f \mapsto \|f\|_{k,p}\) 的完备化。
\item
  设 V 是 \(R^{n}\) 的开子集，\(\varphi\colon U\to V\) 是 \(C^{k}\) 类微分同胚。假设 \(\leqslant k\) 阶偏导数——\(\varphi\) 的（相应地 \(\varphi^{-1}\) 的）——在 U 上（相应地在 V 上）有界且一致连续。映射 \(f\mapsto f\circ\varphi\) 是从 \(W^{k,p}(V)\) 到 \(W^{k,p}(U)\) 的拓扑向量空间同构。（利用上一问题的结果。）
\end{enumerate}

\begin{enumerate}
\def\labelenumi{\arabic{enumi})}
\setcounter{enumi}{11}
\tightlist
\item
  \begin{enumerate}
  \def\labelenumii{\alph{enumii})}
  \tightlist
  \item
    对所有整数 \(n \in \mathbf{N}\) 与 \(k \in \mathbf{N}\) ，以及每个实数 \(p \geqslant 1\) ，我们有 \(\mathrm{W}^{k,p}(\mathbf{R}^n) = \mathrm{W}_0^{k,p}(\mathbf{R}^n)\) 。
  \end{enumerate}
\end{enumerate}

\begin{enumerate}
\def\labelenumi{\alph{enumi})}
\setcounter{enumi}{1}
\tightlist
\item
  设 \(p \geqslant 1\) 是实数。设 \(U \subset R^{2}\) 是由 \((x, y) \in \mathbf{R}^{2}\) （满足 0 \textless{} \textbar x\textbar{} \textless{} 1 与 0 \textless{} y \textless{} 1）组成的集。特征函数 \(\varphi\) （\(U \cap (\mathbf{R}_{+}^{*} \times \mathbf{R})\) 的）属于 \(W^{1,p}(U)\) ；若 \(\varepsilon > 0\) 充分小，则不存在 \(C^{1}\) 类的函数 f （在 \(\overline{U}\) 的邻域中）使得 \(\|f - \varphi\|_{1,p} < \varepsilon\) 。
\end{enumerate}

\begin{enumerate}
\def\labelenumi{\arabic{enumi})}
\setcounter{enumi}{12}
\tightlist
\item
  设 U 是 \(\mathbf{R}^n\) 的开子集，\(j\colon \mathrm{U}\to \mathbf{R}^n\) 是典范单射。对每个函数 \(f\colon \mathrm{U}\to \mathbf{C}\) ，我们用 \(j_{!}f\) 表示 \(f\) 到 \(\mathbf{R}^n\) 的零延拓，即从 \(\mathbf{R}^n\) 到 \(\mathbf{C}\) 的、在 U 上与 \(f\) 一致且在 U 外为零的函数。
\end{enumerate}

设 \(k \geqslant 0\) 是整数，\(p \geqslant 1\) 是实数。

\begin{enumerate}
\def\labelenumi{\alph{enumi})}
\item
  设 \(f \in \mathrm{W}^{k,p}(\mathrm{U})\) 。对每个 \(\alpha \in \mathbf{N}^n\) （满足 \(|\alpha| \leqslant k\) ），我们有 \(\partial^\alpha(j_{!}f) = j_{!}(\partial^\alpha(f))\) 在 \(\mathcal{D}'(\mathbf{R}^n)\) 中。
\item
  由向子空间的过渡诱导的映射 \(f \mapsto j_{!}f\) 是从 \(\mathrm{W}_{0}^{k,p}(\mathrm{U})\) 到 \(\mathrm{W}^{k,p}(\mathbf{R}^{n})\) 中的等距线性映射。
\end{enumerate}

\begin{enumerate}
\def\labelenumi{\arabic{enumi})}
\setcounter{enumi}{13}
\tightlist
\item
  设 U 是 \(R^{n}\) 的开子集，\(k \geqslant 1\) 是整数。设 \(p \geqslant 1\) 是实数。假设余集 \(R^{n} - U\) 关于勒贝格测度不是可忽略的。用 j 表示 U 到 \(R^{n}\) 中的典范单射。
\end{enumerate}

\begin{enumerate}
\def\labelenumi{\alph{enumi})}
\item
  存在 \(R^{n}\) 中相对紧的开长方体 B，使得 \(B \cap U\) 与 \(B \cap (\mathbf{R}^{n} - \mathbf{U})\) 都不是可忽略的。
\item
  设 \(\widetilde{f} \in \mathcal{D}(\mathbf{R}^n)\) 在 \(B \cap U\) 上等于 1。设 \(f\) 是 \(\widetilde{f}\) 到 U 的限制。那么 \(f\) 不属于 \(\mathrm{W}_0^{k,p}(\mathrm{U})\) （否则，零延拓 \(j_{!}f\) 将属于 \(\mathrm{W}^{k,p}(\mathbf{R}^n)\) ；\(j_{!}f\) 到 B 的限制将是一个满足 \(\partial_i(j_{!}f) = 0\) （对 \(1 \leqslant i \leqslant n\) ）的广义函数，从而将是与一个常数相伴的广义函数。）
\end{enumerate}

\begin{enumerate}
\def\labelenumi{\arabic{enumi})}
\setcounter{enumi}{14}
\tightlist
\item
  设 \(n \geqslant 1\) 是整数，p 是满足 \(1 \leqslant p < n\) 的实数。
\end{enumerate}

\begin{enumerate}
\def\labelenumi{\alph{enumi})}
\item
  设 \(\alpha\) 是满足 \(0 < \alpha < n/p - 1\) 的实数。设 \(g\) 是 \(\mathbf{R}_{+}^{*}\) 上的无穷可微正函数，使得 \(g(x) = x^{-\alpha}\) （当 \(x\) 充分小）且 \(g(x) = 0\) （当 \(x\) 充分大）；从 \(\mathbf{R}^{n}\) 到 \(\mathbf{C}\) 的、由 \(x \mapsto g(\|x\|)\) 定义的函数属于 \(\mathrm{W}^{1,p}(\mathbf{R}^{n})\) 。
\item
  设 \(i \mapsto q_{i}\) 是从 N 到 \(Q^{n}\) 的双射，f 是由下式定义的从 \(R^{n}\) 到 C 的映射：
\end{enumerate}

\[
f (x) = \sum_ {i \in \mathbf {N}} 2 ^ {- i} g (\| x - q _ {i} \|)
\]

对 \(x \in \mathbf{R}^n\) 成立。定义 \(f\) 的级数在 \(\mathrm{W}^{1,p}(\mathbf{R}^n)\) 中收敛；对 \(\mathbf{R}^n\) 的每个开集 U，\(f\) 到 U 的限制都不是几乎处处可微的。\(c)\) 存在在任一点都不连续的函数 \(f \in \mathrm{W}^{1,n}(\mathbf{R}^n)\) 。（对函数 \(x \mapsto \log(|\log(|x|)|)\) 用类似的方法。）

\begin{enumerate}
\def\labelenumi{\arabic{enumi})}
\setcounter{enumi}{15}
\tightlist
\item
  设 \(n \geqslant 1\) 与 \(k \geqslant 1\) 是整数。设 \(p \geqslant 1\) 是实数。设 U 是 \(\mathbf{R}^n\) 的有界开子集。
\end{enumerate}

\begin{enumerate}
\def\labelenumi{\alph{enumi})}
\item
  假设 \(n \neq kp\) 。设 q 是满足 \(1 \leqslant q \leqslant np/(n-kp)\) 的实数。\(\mathcal{D}(U)\) 到 \(L^{q}(U)\) 中的单射容许一个连续延拓 \(i\) （从 \(W_{0}^{k,p}(U)\) 到 \(L^{q}(U)\) 中）。
\item
  假设 \(1 \leqslant q \leqslant np/(n-kp)\) 。线性映射 \(i\) 是紧的（“雷利希–孔德拉绍夫定理”）。
\item
  假设 \(k \geqslant 1\) 且 \(q = np/(n - p)\) 。线性映射 i 不是紧的。
\end{enumerate}

\begin{enumerate}
\def\labelenumi{\arabic{enumi})}
\setcounter{enumi}{16}
\tightlist
\item
  设 \(n \geqslant 1\) 是整数。我们用 \(\mu\) 表示 \(\mathbf{R}^n\) 上的勒贝格测度。
\end{enumerate}

\begin{enumerate}
\def\labelenumi{\alph{enumi})}
\tightlist
\item
  \(\mathrm{L}^{2}(\mathbf{R}^{n},\mu)\) 的有界子集 B 是紧的，当且仅当对每个 \(\varepsilon>0\) ，存在 \(R^{n}\) 的紧子集 K 使得
\end{enumerate}

\[
\int_ {\mathbf {R} ^ {n} - K} | f | ^ {2} + \int_ {\mathbf {R} ^ {n} - K} | \widehat {f} | ^ {2} <   \varepsilon
\]

对所有 \(f\in \mathrm{B}\) 成立。

\begin{enumerate}
\def\labelenumi{\alph{enumi})}
\setcounter{enumi}{1}
\tightlist
\item
  设 \(k \in N\) 。设 K 是 \(R^{n}\) 的紧子集。我们用 \(\mathrm{H}_{k}^{2}(\mathrm{K})\) 表示索博列夫空间 \(\mathrm{H}_{k}^{2}(\mathbf{R}^{n})\) 的由在紧集 K 外为零的函数组成的闭子空间。典范单射 \(\mathrm{H}_{k}^{2}(\mathrm{K}) \to \mathrm{H}_{k}^{2}(\mathbf{R}^{n})\) 是紧的。
\end{enumerate}

\begin{enumerate}
\def\labelenumi{\arabic{enumi})}
\setcounter{enumi}{17}
\tightlist
\item
  设 \(n \in \mathbf{N}\) ，\(U\) 是 \(\mathbf{R}^n\) 的开子集。本习题的目标是证明 \(\mathcal{D}'(U)\) 是有界型空间。
\end{enumerate}

若 K 是 U 的紧子集，我们给 \(\mathcal{C}_{\mathrm{K}}^{\infty}(\mathrm{U})'\) 赋予有界收敛拓扑。我们用 \(\pi_{\mathrm{K}}\colon\mathcal{D}'(\mathrm{U})\to\mathcal{C}_{\mathrm{K}}^{\infty}(\mathrm{U})'\) 表示把广义函数映为其到 \(\mathcal{C}_{\mathrm{K}}^{\infty}(\mathrm{U})\) 的限制的线性映射；它是连续的。

设 \(S \subset \mathcal{D}'(U)\) 是凸的、均衡的且吸收有界集的子集。

\begin{enumerate}
\def\labelenumi{\alph{enumi})}
\item
  存在 U 的紧子集 K 使得 \(\mathrm{Ker}(\pi_{\mathrm{K}}) \subset \mathrm{S}\) 。（否则，考虑 U 的紧子集的递增序列 \((\mathrm{K}_{m})_{m \in \mathbf{N}}\) ，其诸内部覆盖 U，并选取 \(f_{m}\) 于 \(\mathrm{Ker}(\pi_{\mathrm{K}_{m}}) - \mathrm{S}\) 中；那么诸元素 \(mf_{m}\) （\(m \in N\) ）所成之集是一个不被 S 吸收的有界子集。）
\item
  由广义函数 \(f \in \mathcal{D}'(\mathrm{U})\) （满足 \(f + \operatorname{Ker}(\pi_{\mathrm{K}}) \subset \mathrm{S}\) ）组成的集 T 是凸的、均衡的且吸收有界集的集。
\item
  集 \(\pi_{\mathrm{K}}(\mathrm{T})\) 是 \(\mathcal{C}_{\mathrm{K}}^{\infty}(\mathrm{U})'\) 的凸的、均衡的且吸收有界集的子集。（设 A 是 \(\mathcal{C}_{\mathrm{K}}^{\infty}(\mathrm{U})'\) 的有界子集；证明存在 \(\mathcal{D}(\mathrm{U})\) 中 0 的邻域 W 使得 \(\mathrm{A} \subset (\mathrm{W} \cap \mathcal{C}_{\mathrm{K}}^{\infty}(\mathrm{U}))^{\circ}\) ，并由此推出 A 含于 \(\pi_{\mathrm{K}}(\mathrm{W}^{\circ})\) ，从而被 \(\pi_{\mathrm{K}}(\mathrm{T})\) 吸收。）
\end{enumerate}

\(d)\) 由此推出 S 是 0 的邻域，并断言空间 \(\mathcal{D}'(U)\) 是有界型的。（注意 \(\pi_{\mathrm{K}}(\mathrm{T})\) 是 0 的邻域，并注意 S 包含 \(\pi_{\mathrm{K}}^{-1}(\pi_{\mathrm{K}}(\mathrm{T}))\) 。）

\begin{enumerate}
\def\labelenumi{\arabic{enumi})}
\setcounter{enumi}{18}
\tightlist
\item
  设 \(n \in \mathbf{N}\) ，\(\mathrm{U} \subset \mathbf{R}^n\) 是开集。
\end{enumerate}

\begin{enumerate}
\def\labelenumi{\alph{enumi})}
\tightlist
\item
  设 \(f \in \mathcal{D}'(\mathrm{U})\) 。开集 \(\mathrm{W} \subset \mathrm{U}\) （使得 \(f\) 到 W 的限制为零）的并 V 在 U 中是开的。
\end{enumerate}

f 的支撑称为 V 的余集；它是 U 的闭子集。

\begin{enumerate}
\def\labelenumi{\alph{enumi})}
\setcounter{enumi}{1}
\item
  设 \(f \in \mathcal{D}'(\mathrm{U})\) ，\(\varphi \in \mathcal{D}(\mathrm{U})\) 的支撑与 f 的支撑不相交。我们有 \(\langle f, \varphi \rangle = 0\) 。
\item
  设 \(f \in \mathcal{D}'(\mathrm{U})\) 且 \(\alpha \in N^{n}\) 。\(\partial^{\alpha}f\) 的支撑含于 f 的支撑。
\item
  若 \(f \in \mathcal{D}'(\mathrm{U})\) 是与测度 \(\nu\) 相伴的广义函数，则 f 的支撑与 \(\nu\) 的支撑重合。
\item
  设 \(f\) 是 U 中具有紧支撑的广义函数。线性形式 \(f\) 可延拓为空间 \(\mathcal{C}^{\infty}(\mathrm{U})\) 上的连续线性形式。
\item
  U 中具有紧支撑的广义函数所成的空间可等同于空间 \(\mathcal{C}^{\infty}(\mathrm{U})\) 的对偶。
\end{enumerate}

\(g)\) 设 \(f \in \mathcal{D}'(\mathbf{R}^n)\) 是具有紧支撑的广义函数。我们有 \(f \in \mathcal{S}'(\mathbf{R}^n)\) 。\(h)\) 设 \(f \in \mathcal{D}'(\mathbf{R}^n)\) 是具有紧支撑的广义函数。\(f\) 的傅里叶变换属于 \(\mathcal{C}^\infty(\mathbf{R}^n)\) ；它是连同其所有导数都具有多项式增长的函数。

\begin{enumerate}
\def\labelenumi{\arabic{enumi})}
\setcounter{enumi}{19}
\tightlist
\item
  设 \(n \in \mathbf{N}\) 。我们用 \(V_n\) 表示 \(\mathbf{C}\) 的由复数 \(\lambda\) （它们满足 \(\mathcal{R}(\lambda) > -n - 1\) 且 \(\lambda \notin \{-1, -2, \ldots, -n\}\) ）组成的开集。
\end{enumerate}

设 \(\lambda\in V_{0}\) 。我们用 \(x_{+}^{\lambda}\) 表示 \(\mathbf{R}\) 上的由局部可积函数

\[
x \mapsto \left\{ \begin{array}{l l} 0 & \text { if } x \leqslant 0, \\ x ^ {\lambda} & \text { if } x > 0. \end{array} \right.
\]

\begin{enumerate}
\def\labelenumi{\alph{enumi})}
\item
  映射 \(\lambda \mapsto x_{+}^{\lambda}\) （从 \(\mathrm{V}_0\) 到 \(\mathcal{D}'(\mathbf{R})\) 中的）是解析的。
\item
  设 \(n \in \mathbf{N}\) 。对每个 \(\lambda \in V_n\) ，在 \(\mathcal{D}(\mathbf{R})\) 上由
\end{enumerate}

\[
\begin{array}{l} \varphi \mapsto \int_ {0} ^ {1} x ^ {\lambda} \Big (\varphi (x) - \varphi (0) - x \varphi^ {\prime} (0) - \dots - \frac {x ^ {n - 1}}{(n - 1) !} \varphi^ {(n - 1)} (0) \Big) d x \\ \qquad + \int_ {1} ^ {+ \infty} x ^ {\lambda} \varphi (x) d x + \sum_ {k = 1} ^ {n} \frac {\varphi^ {(k - 1)} (0)}{(k - 1) ! (\lambda + k)} \end{array}
\]

定义的映射是缓增广义函数；它不依赖于使 \(\lambda \in V_{n}\) 的整数 n，并且与 \(x_{+}^{\lambda}\) 一致当 \(\lambda \in V_{0}\) 时。

我们将用 \(x_{+}^{\lambda}\) 表示这样定义的广义函数，对每个 \(\lambda\) 属于诸集 \(V_{n}\) （\(n \in N\) ）之并。

\begin{enumerate}
\def\labelenumi{\alph{enumi})}
\setcounter{enumi}{2}
\tightlist
\item
  设 \(n \in N\) ，\(\lambda \in C\) 满足 \(-n - 1 < \mathcal{R}(\lambda) < -n\)。对 R 上每个施瓦茨函数 \(\varphi\) ，我们有
\end{enumerate}

\[
\langle x _ {+} ^ {\lambda}, \varphi \rangle = \int_ {0} ^ {+ \infty} x ^ {\lambda} \left(\varphi (x) - \varphi (0) - x \varphi^ {\prime} (0) - \dots - \frac {x ^ {n - 1}}{(n - 1) !} \varphi^ {(n - 1)} (0)\right) d x.
\]

  从 \(V_{n}\) 到 \(\mathcal{D}'(\mathrm{U})\) 中的由 \(\lambda\mapsto x_{+}^{\lambda}\) 定义的映射是解析的。

\(d)\) 对每个 \(\lambda\) （属于诸集 \(V_{m}\) 之并的），我们有公式

\[
(x _ {+} ^ {\lambda}) ^ {\prime} = \lambda x _ {+} ^ {\lambda - 1}.
\]

\begin{enumerate}
\def\labelenumi{\arabic{enumi})}
\setcounter{enumi}{20}
\tightlist
\item
  \begin{enumerate}
  \def\labelenumii{\alph{enumii})}
  \tightlist
  \item
    设 \(f \in \mathcal{S}(\mathbf{R})\) 。R 上满足 \(g(0) = 2f'(0)\) 且 \(g(x) = (f(x) - f(-x))/x\) （对所有 \(x \neq 0\) ）的函数 g 在 R 上关于勒贝格测度可积。把 f 映为 g 在 R 上的积分的映射是一个缓增广义函数，记作 vp，称为 1/x 的柯西主值。
  \end{enumerate}
\end{enumerate}

\begin{enumerate}
\def\labelenumi{\alph{enumi})}
\setcounter{enumi}{1}
\tightlist
\item
  对每个 \(f \in \mathcal{S}(\mathbf{R})\) ，我们有
\end{enumerate}

\[
\langle \mathrm{vp}, f \rangle = \lim _ {\varepsilon \rightarrow 0} \int_ {\mathbf {R} - [ - \varepsilon , \varepsilon ]} \frac {f (x)}{x} d x.
\]

\begin{enumerate}
\def\labelenumi{\alph{enumi})}
\setcounter{enumi}{2}
\item
  vp 到 \(\mathbf{R} - \{0\}\) 的限制是与局部可积函数 \(x \mapsto 1/x\) （\(\mathbf{R} - \{0\}\) 上的）相伴的广义函数。
\item
  广义函数 vp 与下述广义函数的导数重合：该广义函数与 R 上在 0 处取值 0 且对 \(x \neq 0\) 由 \(x \mapsto \log(|x|)\) 定义的局部可积函数相伴。
\item
  设 \(\varepsilon > 0\) 。我们用 \(p_{\varepsilon}\) 表示函数 \(x \mapsto 1 / (x + i\varepsilon)\) （\(\mathbf{R}\) 上的）。我们有
\end{enumerate}

\[
\lim _ {\varepsilon \to 0} p _ {\varepsilon} = \mathrm{vp} - i \pi \varepsilon_ {0}
\]

在赋有弱拓扑的空间 \(\mathcal{S}'(\mathbf{R})\) 中成立，其中 \(\varepsilon_0\) 是与 0 处的点质量相伴的缓增广义函数。

\begin{enumerate}
\def\labelenumi{\arabic{enumi})}
\setcounter{enumi}{21}
\item
  R 上与指数函数相伴的广义函数不是缓增的。
\item
  设 \(n \in \mathbf{N}\) ，U 是 \(\mathbf{R}^n\) 的开集。我们用 \(\mu\) 表示 U 上的勒贝格测度。对 \(m \in \mathbf{N}\) ，我们记 \(\boldsymbol{m} = (m, \ldots, m) \in \mathbf{N}^n\) 。
\end{enumerate}

\begin{enumerate}
\def\labelenumi{\alph{enumi})}
\tightlist
\item
  设 \(f \in \mathcal{D}(\mathrm{U})\) 。假设存在整数 \(m \geqslant 1\) 以及实数 \(c \geqslant 0\) 使得
\end{enumerate}

\[
| \langle f, \varphi \rangle | \leqslant c \int_ {\mathrm{U}} | \partial^ {\boldsymbol {m}} \varphi | d \mu
\]

对所有 \(\varphi\in\mathcal{D}(\mathrm{U})\) 成立。于是存在函数 \(g\in\mathrm{L}^{\infty}(\mathrm{U})\) 使得 \(f=\partial^{m}g\) 。（线性映射 \(\varphi\mapsto\partial^{m}\varphi\) 在 \(\mathcal{D}(\mathrm{U})\) 上是单射；设 E 是它的像；考虑 E 上把 \(\partial^{m}\varphi\) 映为 \(\langle f,\varphi\rangle\) 的到 C 中的线性形式，并应用哈恩–巴拿赫定理。）

\begin{enumerate}
\def\labelenumi{\alph{enumi})}
\setcounter{enumi}{1}
\tightlist
\item
  设 V 是含于 U 的相对紧开集。存在整数 \(m \geqslant 1\) 以及 \(c \geqslant 0\) 使得
\end{enumerate}

\[
| \langle f, \varphi \rangle | \leqslant c \int_ {\mathrm{U}} | \partial^ {\boldsymbol {m}} \varphi | d \mu
\]

对所有支撑含于 V 的 \(\varphi\in\mathcal{D}(\mathrm{U})\) 成立。

\begin{enumerate}
\def\labelenumi{\alph{enumi})}
\setcounter{enumi}{2}
\tightlist
\item
  对每个相对紧开集 \(V \subset U\) ，存在连续函数 \(g \in \mathcal{C}(V)\) 以及 \(\alpha \in N^{n}\) ，使得 f 到 V 的限制等于 \(\partial^{\alpha}g\)。\(d)\) 给出一个广义函数 \(f \in \mathcal{D}'(\mathbf{R}^{n})\) 的例子，使得 f 不是 \(f = \partial^{\alpha}g\) 的形式，其中 \(g \in \mathcal{C}(\mathbf{R}^{n})\) 是连续函数且 \(\alpha \in N^{n}\) 。
\end{enumerate}

\begin{enumerate}
\def\labelenumi{\arabic{enumi})}
\setcounter{enumi}{23}
\tightlist
\item
  设 \(n \in N\) ，U 是 \(R^{n}\) 的开集。设 \(f \in \mathcal{D}'(U)\) 是支撑 C 在 U 中为紧的广义函数。
\end{enumerate}

\begin{enumerate}
\def\labelenumi{\alph{enumi})}
\tightlist
\item
  存在整数 \(m \in N\) ，使得广义函数 f 延拓为空间 \(\mathcal{C}^{m}(U)\) 上的连续线性形式，该空间赋有函数及其阶 \(\leqslant m\) 的偏导数在紧子集上一致收敛的拓扑。
\end{enumerate}

满足该性质的最小整数 m 称为 f 的阶。

\begin{enumerate}
\def\labelenumi{\alph{enumi})}
\setcounter{enumi}{1}
\item
  设 \(\varphi \in \mathcal{D}(\mathrm{U})\) 使得 \(\partial^{\alpha}\varphi\) 对每个 \(\alpha \in \mathbf{N}^{n}\) （满足 \(|\alpha| \leqslant m\) ）在 C 上为零。我们则有 \(\langle f, \varphi \rangle = 0\) 。（对充分小的 \(\delta > 0\) ，考虑 C 的邻域 \(\mathrm{U}_{\delta}\) ，它由 \(x \in \mathrm{U}\) （满足 \(d(x, \mathrm{C}) \leqslant \delta\) ）组成，并证明存在 \(\eta \geqslant 0\) （当 \(\delta\) 趋于 0 时趋于 0），使得对某个 \(\mathrm{C}_{0} \geqslant 0\) ，有 \(|\partial^{\alpha}\varphi(x)| \leqslant \mathrm{C}_{0}\delta^{m-|\alpha|}\eta\) 对 \(x \in \mathrm{U}_{\delta}\) 与 \(|\alpha| \leqslant m\) ；证明存在函数 \(\beta_{\delta} \in \mathcal{D}(\mathrm{U})\) ，它在 \(\mathrm{U}_{\delta/4}\) 上等于 1 且支撑含于 \(\mathrm{U}_{\delta}\) ，并满足 \(|\partial^{\alpha}\beta_{\delta}| \leqslant \mathrm{C}_{1}\delta^{-|\alpha|}\) 对某个 \(\mathrm{C}_{1} \geqslant 0\) ；最后注意 \(\langle f, \varphi \rangle = \langle f, \varphi\beta_{\delta} \rangle\) 且 \(\varphi\beta_{\delta}\) 在 \(\mathcal{C}^{m}(\mathrm{U})\) 中趋于 0。）
\item
  设 \(x \in U\) 。支撑化为单点 x 的广义函数 f 所成的空间以诸广义函数 \(\partial^{\alpha}\varepsilon_{x}\) （其中 \(\alpha \in N^{n}\) ）为代数基。
\item
  设 \(n = 1\) 。公式
\end{enumerate}

\[
\langle f, \varphi \rangle = \lim _ {m \rightarrow + \infty} \left(\sum_ {1 \leqslant k \leqslant m} \varphi (k ^ {- 1}) - m \varphi (0) - \log (m) \varphi^ {\prime} (0)\right)
\]

在 \(\mathbf{R}\) 上定义了一个阶为 1 的广义函数，其支撑是紧的，等于集 \(C = \{0\} \cup \{1/(n+1) \mid n \in \mathbf{N}\}\) 。存在序列 \((\varphi_n)\) （在 \(\mathcal{D}(\mathbf{R})\) 中），使得 \(\partial^{\alpha}\varphi_{n}\) 对每个 \(\alpha\in\mathbf{N}^{n}\) （满足 \(|\alpha|\leqslant1\) ）在 C 上一致收敛于 0，但 \(\langle f,\varphi_{n}\rangle\to+\infty\) 。

\begin{enumerate}
\def\labelenumi{\arabic{enumi})}
\setcounter{enumi}{24}
\tightlist
\item
  设 \(n \in \mathbf{N}\) 。我们用 \(N\) 表示 \(\mathbf{R}^n\) 上的欧几里得范数，并用 \(\mu\) 表示勒贝格测度。
\end{enumerate}

我们用 \(\mathcal{D}_{\mathrm{L}^{1}}(\mathbf{R}^{n})\) 表示 \(\mathcal{C}^{\infty}(\mathbf{R}^{n})\) 的向量子空间，其元素是函数 \(\varphi\) ，它满足 \(\partial^{\alpha}\varphi\in\mathrm{L}^{1}(\mathbf{R}^{n})\) 且 \(\partial^{\alpha}\varphi\in\mathcal{C}_{0}(\mathbf{R}^{n})\) 对所有 \(\alpha\in\mathbf{N}^{n}\) 。\(a)\) 空间 \(\mathcal{D}_{\mathrm{L}^{1}}(\mathbf{R}^{n})\) 赋有由半范数

\[
q _ {\alpha} (\varphi) = \int_ {\mathbf {R} ^ {n}} | \partial^ {\alpha} \varphi | d \mu , \quad \alpha \in \mathbf {N} ^ {n},
\]

定义的拓扑后是弗雷歇空间。

\begin{enumerate}
\def\labelenumi{\alph{enumi})}
\setcounter{enumi}{1}
\item
  空间 \(\mathcal{D}(\mathbf{R}^n)\) 含于 \(\mathcal{D}_{\mathrm{L}^1}(\mathbf{R}^n)\) 且在其中稠密。可以把对偶空间 \(\mathcal{D}_{\mathrm{L}^1}'(\mathbf{R}^n)\) （\(\mathcal{D}_{\mathrm{L}^1}(\mathbf{R}^n)\) 的）等同于 \(\mathcal{D}'(\mathbf{R}^n)\) 的一个子空间。
\item
  设 \(f \in \mathcal{D}'(\mathbf{R}^n)\) 。我们有 \(f \in \mathcal{D}_{\mathrm{L}^1}'(\mathbf{R}^n)\) 当且仅当存在有限集 I、族 \((\alpha_i)_{i \in \mathrm{I}}\) （在 \(\mathbf{N}^n\) 中）以及族 \((g_i)_{i \in \mathrm{I}}\) （在 \(\mathrm{L}^\infty(\mathbf{R}^n)\) 中）使得
\end{enumerate}

\[
f = \sum_ {i \in I} \partial^ {\alpha_ {i}} g _ {i}.
\]

（确定 I 与族 \((\alpha_{i})\) ，使得 \(|\langle f,\varphi\rangle |\) 有界，当 \(\varphi\) 满足 \(q_{\alpha_i}(\varphi)\leqslant 1\) 对所有 \(i\in \mathrm{I}\) ；把 \(\mathcal{D}_{\mathrm{L}^1}(\mathbf{R}^n)\) 等同于 \(\mathrm{L}^1 (\mathbf{R}^n)^m\) 的一个向量子空间，借助映射 \(\varphi \mapsto (\partial^{\alpha_i}\varphi)_{i\in \mathrm{I}}\) ，并应用哈恩–巴拿赫定理。）

\begin{enumerate}
\def\labelenumi{\alph{enumi})}
\setcounter{enumi}{3}
\tightlist
\item
  设 \(f \in \mathcal{D}'(\mathbf{R}^n)\) 。我们有 \(f \in \mathcal{D}_{\mathrm{L}^1}'(\mathbf{R}^n)\) 当且仅当存在具有缓和增长的连续函数 \(g\) （在 \(\mathbf{R}^n\) 上）以及 \(\alpha \in \mathbf{N}^n\) 使得 \(f = \partial^\alpha g\)。\(e)\) 设 \(f \in \mathcal{D}'(\mathbf{R}^n)\) 。存在整数 \(m \in \mathbf{N}\) 使得 \((1 + \mathrm{N}^2)^{-m/2}f\) 属于 \(\mathcal{D}_{\mathrm{L}^1}'(\mathbf{R}^n)\) 。
\end{enumerate}

\(f)\) 设 \(f \in \mathcal{D}'(\mathbf{R}^n)\) 。我们有 \(f \in \mathcal{S}'(\mathbf{R}^n)\) 当且仅当存在有界连续函数 \(g\colon \mathbf{R}^n \to \mathbf{C}\) 、整数 \(m \in \mathbf{N}\) 以及 \(\alpha \in \mathbf{N}^n\) 使得 \(f = \partial^\alpha((1 + \mathrm{N}^2)^{m/2}g)\) 。

\begin{enumerate}
\def\labelenumi{\arabic{enumi})}
\setcounter{enumi}{25}
\tightlist
\item
  设 \(n \in N\) 。我们给 \(R^{n}\) 赋予勒贝格测度 \(\mu\) ，并把 \(R^{n}\) 与其对偶通过映射 \((x, y) \mapsto \exp(2i\pi x \cdot y)\) 等同。
\end{enumerate}

\begin{enumerate}
\def\labelenumi{\alph{enumi})}
\item
  对任何函数 \(f \in \mathcal{S}(\mathbf{R}^n)\) ，映射 \(g \mapsto f * g\) 是局部凸空间 \(\mathcal{S}(\mathbf{R}^n)\) 的自同构。
\item
  设 \(f \in \mathcal{S}'(\mathbf{R}^{n})\) 且 \(g \in \mathcal{S}(\mathbf{R}^{n})\) 。映射 \(\varphi \mapsto \langle f, \check{g} * \varphi \rangle\) 是一个缓增广义函数，称为 f 与 g 的卷积。若 \(f \in \mathcal{S}(\mathbf{R}^{n})\) ，它与同 \(f * g\) 相伴的广义函数一致。我们仍用 \(f * g\) 表示这样定义的缓增广义函数。
\end{enumerate}

在下面的问题中，我们固定 \(f \in \mathcal{S}'(\mathbf{R}^n)\) 与 \(g \in \mathcal{S}(\mathbf{R}^n)\) 。\(c)\) 我们有 \(f * g \in \mathcal{C}^\infty(\mathbf{R})\) 并且

\[
\partial^ {\alpha} (f * g) = \partial^ {\alpha} (f) * g = f * \partial^ {\alpha} (g)
\]

对所有 \(\alpha\in\mathbf{N}^{n}\) 成立。

\(d)\) 对每个 \(y \in \mathbf{R}^n\)，用 \(g_y\) 表示函数 \(x \mapsto g(y - x)\) （\(\mathbf{R}^n\) 上的）。函数 \(f * g\) 满足

\[
(f * g) (y) = \langle f, g _ {y} \rangle
\]

对所有 \(y \in R^{n}\) 成立。

\begin{enumerate}
\def\labelenumi{\alph{enumi})}
\setcounter{enumi}{4}
\item
  设 \(h \in \mathcal{S}(\mathbf{R}^n)\) 。我们有 \((f*g)*h = f*(g*h)\) 。
\item
  函数 \(f*g\) 具有多项式增长，其所有导数亦然。（利用习题 25。）
\item
  我们有 \(\mathcal{F}(f*g) = \mathcal{F}(f)\mathcal{F}(g)\) 。
\item
  fg 的傅里叶变换是 \(R^{n}\) 上的无穷可微函数，并且连同其所有导数都具有多项式增长。此外，对每个 \(y \in R^{n}\) ，我们有
\end{enumerate}

\[
\mathcal {F} (f g) (y) = \langle f, g e _ {- y} \rangle
\]

其中 \(e_y\) 表示特征标 \(x \mapsto \exp(2i\pi x \cdot y)\) （\(\mathbf{R}^n\) 的）。

\begin{enumerate}
\def\labelenumi{\arabic{enumi})}
\setcounter{enumi}{26}
\tightlist
\item
  设 \(n \in \mathbf{N}\)。我们用 N 表示 \(\mathbf{R}^{n}\) 上的欧几里得范数，并用 \(\mu\) 表示 \(\mathbf{R}^{n}\) 的每个开子集上的勒贝格测度。设 \(\Delta\) 是微分算子
\end{enumerate}

\[
\Delta = - \sum_ {i = 1} ^ {n} \partial_ {i} ^ {2}
\]

将它视为 \(\mathcal{D}'(\mathbf{R}^{n})\) 的自同态。

对 \(k \in \mathbf{R}\) ，我们用 \(\mathrm{S}^k(\mathbf{R}^n)\) 表示这样的缓增广义函数 \(f\) （\(\mathcal{S}'(\mathbf{R}^n)\) 中的）所成的空间：\((1 + \mathrm{N}^k)\mathcal{F}(f)\) 属于 \(\mathrm{L}^2(\mathbf{R}^n)\) 。

设 \(U \subset R^{n}\) 是 \(R^{n}\) 的开子集。我们用 \(\mathrm{S}_{\mathrm{U}}^{k}(\mathbf{R}^{n})\) 表示满足如下条件的广义函数 \(f \in \mathcal{D}'(\mathbf{R}^{n})\) 所成的空间：对每个支撑含于 U 的试验函数 \(\varphi \in \mathcal{D}(\mathbf{R}^{n})\) ，我们有 \(\varphi f \in \mathrm{S}^{k}(\mathbf{R}^{n})\) 。

\begin{enumerate}
\def\labelenumi{\alph{enumi})}
\tightlist
\item
  空间 \(\mathrm{S}^{k}(\mathbf{R}^{n})\) 赋有内积
\end{enumerate}

\[
\langle f _ {1} \mid f _ {2} \rangle = \int_ {\mathbf {R} ^ {n}} (1 + \mathrm{N} ^ {k}) \overline {{\mathcal {F} (f _ {1})}} \mathcal {F} (f _ {2}) d \mu .
\]

当 \(k \in \mathbf{N}\) 时，我们有 \(W^{k}(\mathbf{R}^{n}) = S^{k}(\mathbf{R}^{n})\) ，并且这两个巴拿赫空间的范数等价；此外，这时我们有 \(W^{k}(\mathbf{R}^{n}) \subset S_{\mathrm{U}}^{k}(\mathbf{R}^{n})\) 。

\begin{enumerate}
\def\labelenumi{\alph{enumi})}
\setcounter{enumi}{1}
\tightlist
\item
  \(\mathrm{S}^{k}(\mathbf{R}^{n})\) 的对偶通过双线性形式等同于 \(\mathrm{S}^{-k}(\mathbf{R}^{n})\)
\end{enumerate}

\[
(f, g) \mapsto \int_ {\mathbf {R} ^ {n}} \mathcal {F} (f) (- x) \mathcal {F} (g) (x) d \mu (x)
\]

（对 \((f,g)\in \mathrm{S}^{-k}(\mathbf{R}^n)\times \mathrm{S}^k (\mathbf{R}^n).\)）

\begin{enumerate}
\def\labelenumi{\alph{enumi})}
\setcounter{enumi}{2}
\tightlist
\item
  假设 U 是相对紧的。我们有
\end{enumerate}

\[
\bigcup_ {k \in \mathbf {R}} \mathrm{S} _ {\mathrm{U}} ^ {k} (\mathbf {R} ^ {n}) = \mathcal {D} ^ {\prime} (\mathbf {R} ^ {n}).
\]

（设 \(f \in \mathcal{D}'(\mathbf{R}^n)\) ；把习题 19 的结果应用于 \(\varphi f\) ，其中 \(\varphi \in \mathcal{D}(\mathbf{R}^n)\) 在 U 上等于 1。）

\(d)\) 设 \(k \in \mathbf{R}\) 且 \(f \in \mathrm{S}^k(\mathbf{R}^n)\) 。对每个整数 \(i\) （满足 \(1 \leqslant i \leqslant n\) ），我们有 \(\partial_i f \in \mathrm{S}^{k-1}(\mathbf{R}^n)\) 。此外，若 \(\Delta(f) \in \mathrm{S}^k(\mathbf{R}^n)\) ，则 \(f \in \mathrm{S}^{k+2}(\mathbf{R}^n)\) 。

\begin{enumerate}
\def\labelenumi{\alph{enumi})}
\setcounter{enumi}{4}
\tightlist
\item
  设 \(k \in \mathbf{R}\) 且 \(f \in \mathrm{S}_{\mathrm{U}}^{k}(\mathbf{R}^{n})\) 。若 \(\Delta(f) \in \mathrm{S}_{\mathrm{U}}^{k}(\mathbf{R}^{n})\) ，则 \(f \in \mathrm{S}_{\mathrm{U}}^{k+2}(\mathbf{R}^{n})\) 。
\end{enumerate}

\(f)\) 设 \(k > n/2\) 是实数，\(m \in \mathbf{N}\) 是满足 \(m < k - n/2\) 的整数。空间 \(\mathrm{S}^k(\mathbf{R}^n)\) 含于 \(\mathcal{C}^m(\mathbf{R}^n)\) （利用柯西–施瓦茨不等式，验证 \(x \mapsto x^\alpha \mathscr{F} f(x)\) 属于 \(\mathrm{L}^1(\mathbf{R}^n)\) 对所有 \(\alpha \in \mathbf{N}^n\) （满足 \(|\alpha| \leqslant m\) ）。然后利用傅里叶反演公式证明 \(\mathscr{F}(f)\) 在 \(\mathbf{R}\) 上可微。）

\(g)\) 设 \(k > n/2\) 是实数，\(m \in \mathbf{N}\) 是满足 \(m < k - n/2\) 的整数。\(\mathrm{S}_{\mathrm{U}}^{k}(\mathbf{R}^{n})\) 的元素到 U 的限制属于 \(\mathcal{C}^{m}(\mathrm{U})\) （对每个 \(x \in \mathrm{U}\) ，验证得到一个良定义的函数 \(g \in \mathcal{C}^{m}(\mathrm{U})\) ，它由 \(g(x) = (\varphi_{x}f)(x)\) 定义，其中 \(\varphi_{x} \in \mathcal{D}(\mathbf{R}^{n})\) 是支撑在 U 中且在 \(x\) 的一个邻域中等于 1 的函数；最后验证 \(f\) 到 U 的限制等于 \(g\) 。）

\begin{enumerate}
\def\labelenumi{\alph{enumi})}
\setcounter{enumi}{7}
\tightlist
\item
  设 \(g \in \mathcal{D}'(\mathbf{R}^n)\) 。若 \(f \in \mathcal{D}'(\mathbf{R}^n)\) 满足 \(\Delta(f) = g\) 且 \(g\) 到 U 的限制属于 \(\mathcal{C}^\infty(\mathrm{U})\) ，则 \(f\) 到 U 的限制属于 \(\mathcal{C}^\infty(\mathrm{U})\) （“椭圆正则性定理”）。
\end{enumerate}

\begin{enumerate}
\def\labelenumi{\arabic{enumi})}
\setcounter{enumi}{27}
\tightlist
\item
  设 \(n\) 与 \(m\) 是整数，\(\mathrm{U} \subset \mathbf{R}^n\) 与 \(\mathrm{V} \subset \mathbf{R}^m\) 是开集。\(a)\) 存在唯一的单射线性映射 \(u\) ，它从 \(\mathcal{D}(\mathrm{U}) \otimes \mathcal{D}(\mathrm{V})\) 到 \(\mathcal{D}(\mathrm{U} \times \mathrm{V})\) 中，使得 \(u(\varphi \otimes \psi)\) 是函数 \((x, y) \mapsto \varphi(x)\psi(y)\) 对所有 \((\varphi, \psi) \in \mathcal{D}(\mathrm{U}) \times \mathcal{D}(\mathrm{V})\) 。此外，\(u\) 的像在 \(\mathcal{D}(\mathrm{U} \times \mathrm{V})\) 中稠密。我们把空间 \(\mathcal{D}(\mathrm{U}) \otimes \mathcal{D}(\mathrm{V})\) 等同于 \(\mathcal{D}(\mathrm{U} \times \mathrm{V})\) 的一个子空间，借助映射 \(u\) 。
\end{enumerate}

\begin{enumerate}
\def\labelenumi{\alph{enumi})}
\setcounter{enumi}{1}
\item
  对 U 上的所有广义函数 \(f\) 与 V 上的所有广义函数 \(g\) ，存在 U × V 上唯一的广义函数 \(f \otimes g\) ，使得 \(\langle f \otimes g, \varphi \otimes \psi \rangle = \langle f, \varphi \rangle \langle g, \psi \rangle\) 对所有 \((\varphi, \psi) \in \mathcal{D}(\mathrm{U}) \times \mathcal{D}(\mathrm{V})\) 成立。
\item
  设 \(f \in \mathcal{D}(\mathrm{U})\) 且 \(g \in \mathcal{D}(\mathrm{V})\) 。对每个 \(\eta \in \mathcal{D}(\mathrm{U} \times \mathrm{V})\) 以及每个 \(x \in U\)，用 \(\eta_x\) 表示映射 \(y \mapsto \eta(x, y)\)。我们有 \(\eta_x \in \mathcal{D}(\mathrm{V})\) ，并且映射 \(\varphi\) （它把每个 \(x \in U\) 映为 \(\langle f, \eta_x \rangle\) ）属于 \(\mathcal{D}(\mathrm{U})\) 且满足
\end{enumerate}

\[
\langle f \otimes g, \eta \rangle = \langle f, \varphi \rangle .
\]

\begin{enumerate}
\def\labelenumi{\arabic{enumi})}
\setcounter{enumi}{28}
\tightlist
\item
  设 \(n\) 与 \(m\) 是整数，\(\mathrm{U} \subset \mathbf{R}^n\) 与 \(\mathrm{V} \subset \mathbf{R}^m\) 是开集。我们沿用上一个习题的记号。设 \(u\) 是从 \(\mathcal{D}(\mathrm{U})\) 到 \(\mathcal{D}'(\mathrm{V})\) 中的连续线性映射。
\end{enumerate}

\begin{enumerate}
\def\labelenumi{\alph{enumi})}
\item
  存在唯一的连续线性形式 \(\widetilde{\kappa}\) ，它定义在子空间 \(\mathcal{D}(\mathrm{U})\otimes \mathcal{D}(\mathrm{V})\) （\(\mathcal{D}(\mathrm{U}\times \mathrm{V})\) 的）上，使得 \(\widetilde{\kappa} (\varphi \otimes \psi) = \langle u(\varphi),\psi \rangle\) 对所有 \((\varphi ,\psi)\) （\(\mathcal{D}(\mathrm{U})\times \mathcal{D}(\mathrm{V})\) 中的）成立。
\item
  存在广义函数 \(\kappa(u)\) （在 \(U \times V\) 上）延拓 \(\widetilde{\kappa}\) ，并且映射 \(u \mapsto \kappa(u)\) （从 \(\mathcal{L}(\mathcal{D}(U), \mathcal{D}'(V))\) 到 \(\mathcal{D}'(U \times V)\) 中的）是连续的。
\item
  映射 \(\kappa\) 是双射（“施瓦茨核定理”；为证明 \(\kappa\) 的满射性，给定 U × V 上的广义函数 \(\eta\) ，把 \(u(\varphi)\) 定义为 V 上的广义函数 \(\psi \mapsto \langle \eta, \varphi \otimes \psi \rangle\) 。）
\end{enumerate}

\begin{enumerate}
\def\labelenumi{\arabic{enumi})}
\setcounter{enumi}{29}
\tightlist
\item
  设 \(n \in \mathbf{N}\) 。我们用 \((x_1, \ldots, x_n)\) 表示 \(\mathbf{R}^n\) 的典范基的对偶基。
\end{enumerate}

\begin{enumerate}
\def\labelenumi{\alph{enumi})}
\tightlist
\item
  设 \(c \geqslant 0\) 与 r \textgreater{} 0 是实数。存在在 \(R^{n}\) 中 0 的开邻域 U 上定义且解析的函数 u 使得
\end{enumerate}

\[
\left(r - \sum_ {j = 1} ^ {n - 1} x _ {j} - u\right) \partial_ {n} u - c r \sum_ {j = 1} ^ {n - 1} \partial_ {j} u = c r
\]

在 U 上成立，并且 \(u(x,0)=0\) 对 \(x\in\mathbf{R}^{n-1}\) （满足 \((x,0)\in\mathrm{U}\) ）。（先考虑 n=2 的情形。）

\begin{enumerate}
\def\labelenumi{\alph{enumi})}
\setcounter{enumi}{1}
\tightlist
\item
  设 U 是 \(\mathbf{R}^{n}\) 中 0 的开邻域。设 \((a_{1},\ldots,a_{n-1},b)\) 是从 \(\mathrm{U}\times\mathbf{R}\) 到 \(\mathbf{C}\) 的一组函数，它们在 \((0,0)\in\mathrm{U}\times\mathbf{R}\) 的一个邻域中解析。我们用 Q 表示从 U 上解析函数所成的向量空间到其自身的映射，使得
\end{enumerate}

\[
\mathrm{Q} (u) = \partial_ {n} u - \sum_ {j = 1} ^ {n - 1} a _ {j} (x _ {1}, \dots , x _ {n - 1}, u) \partial_ {j} u + b (x _ {1}, \dots , x _ {n - 1}, u).
\]

至多存在一个函数 \(v\) ，它在 0 的开邻域 \(V\) （含于 \(U\) ）中定义且解析，使得 \(Q(v) = 0\) 并且 \(v(x,0) = 0\) 对所有 \(x \in \mathbf{R}^{n-1}\) （满足 \((x,0) \in V\) ）。（验证这样的解 \(v\) 在 0 处的泰勒展开的系数被唯一确定。）

\begin{enumerate}
\def\labelenumi{\alph{enumi})}
\setcounter{enumi}{2}
\tightlist
\item
  设 \(f\) 与 \(g\) 是 \(\mathbf{C}[[z_1, \ldots, z_n]]\) 中的形式幂级数。称 \(f\) 优控 \(g\) ，并记作 \(f \ll g\) ，如果
\end{enumerate}

\[
f = \sum_ {\alpha \in \mathbf {N} ^ {n}} a _ {\alpha} z ^ {\alpha}, \qquad g = \sum_ {\alpha \in \mathbf {N} ^ {n}} b _ {\alpha} z ^ {\alpha}
\]

其中 \(|a_{\alpha}| \leqslant b_{\alpha}\) 对所有 \(\alpha \in \mathbf{N}^{n}\) 。若形式级数 g 在 \(R^{n}\) 中 0 的一个邻域中收敛，则 f 亦然。

\(d)\) 若存在在 0 的邻域中解析的一组函数 \((\widetilde{a}_1, \ldots, \widetilde{a}_{n-1}, \widetilde{b})\) ，使得 \(a_j \ll \widetilde{a}_j\) 对 \(1 \leqslant j \leqslant n-1\) 且 \(b \ll \widetilde{b}\) （即在 0 处的相应泰勒展开满足这些条件的意义下），并且若方程

\[
\partial_ {n} v - \sum_ {j = 1} ^ {n - 1} \widetilde {a} _ {j} (x _ {1}, \dots , x _ {n - 1}, v) \partial_ {j} v + \widetilde {b} (x _ {1}, \dots , x _ {n - 1}, v) = 0
\]

关于未知函数 \(v\) 在 0 的一个邻域中有解析解，则方程 \(\mathrm{Q}(v) = 0\) 在 0 的一个邻域中有唯一的解析解。

\begin{enumerate}
\def\labelenumi{\alph{enumi})}
\setcounter{enumi}{4}
\tightlist
\item
  方程 \(Q(v) = 0\) 存在唯一的在 0 的邻域中解析的解（“柯西–柯瓦列夫斯卡娅定理”）。
\end{enumerate}

\begin{enumerate}
\def\labelenumi{\arabic{enumi})}
\setcounter{enumi}{30}
\tightlist
\item
  我们用 L 表示微分算子
\end{enumerate}

\[
\mathrm{L} = - i \partial_ {1} + \partial_ {2} - 2 (x _ {1} + i x _ {2}) \partial_ {3}
\]

它是 \(\mathbf{R}^{3}\) 上阶为 1 的算子，其中 \((x_{1}, x_{2}, x_{3})\) 是 \(\mathbf{R}^{3}\) 的典范基的对偶基，诸 \(x_{i}\) 被视为 \(\mathbf{R}^{3}\) 上的多项式（“卢伊算子”）。

\begin{enumerate}
\def\labelenumi{\alph{enumi})}
\item
  设 U 是 \(\mathbf{R}^3\) 中 0 的开邻域。存在函数 \(\varphi \in \mathcal{D}(\mathrm{U})\) 使得 \(\varphi(0) = 1\) 并且 \(\mathrm{L}(\varphi)\) 在含于 U 的 0 的邻域 V 中为零。（利用柯西–柯瓦列夫斯卡娅定理。）
\item
  设 \(u \in \mathcal{C}^{\infty}(\mathrm{U})\) 是满足 \(\mathrm{L}(u) = 0\) 且 \(u(0) = 0\) 的函数。假设存在 \(\delta > 0\) 使得 \(\mathcal{R}(u(x)) < -2\delta\) 对 \(x \in U - V\) 成立。我们令 \(v_{t} = \varphi \exp(t(u + \delta))\)。我们有 \(v_{t} \in \mathcal{C}^{\infty}(\mathrm{U})\) ，并且 \(\mathrm{L}(v_{t})\) 在 \(\mathcal{D}(\mathrm{U})\) 中趋于 0 当 \(t \to +\infty\) 时。
\item
  设 F 是由这样的函数 \(f \in \mathcal{C}^{\infty}(U)\) 组成的空间：\(\partial^{\alpha}f\) 对每个 \(\alpha \in N^{3}\) 在 U 上有界。空间 F 赋有半范数 \(q_{\alpha}(f) = \sup_{x \in U} |\partial^{\alpha}f|\) 后是弗雷歇空间。
\end{enumerate}

\(d)\) 设 \(\psi\in\mathcal{D}(\mathbf{R}^{3})\) 。用 \(f_{t}\) 表示函数 \(x\mapsto t^{3}\psi(tx)\) 到 U 的限制。我们有 \(f_{t}\in\mathrm{F}\) 对 \(t\geqslant1\) ，并且存在实数 \(c\geqslant0\) ，使得 \(q_{\alpha}(f_{t})\leqslant ct^{|\alpha|+3}\) 对所有 \(\alpha\in\mathbf{N}^{3}\) 以及每个 \(t\geqslant1\) 。

\begin{enumerate}
\def\labelenumi{\alph{enumi})}
\setcounter{enumi}{4}
\tightlist
\item
  对每个 \(t \geqslant 1\) ，我们用 \(\lambda_{t}\) 表示 F 上的线性形式 \(f \mapsto \int_{U} fv_{t}\)。它是连续的，并且对充分大的 t，我们有
\end{enumerate}

\[
\lambda_ {t} (f _ {t}) = e ^ {\delta t} \left(\int_ {\mathrm{U}} \psi (x) \exp \left(\sum_ {j = 1} ^ {3} \xi_ {j} x _ {j}\right) d x _ {1} d x _ {2} d x _ {3} + \varepsilon (t)\right)
\]

其中 \(\xi_{j} = \partial_{j}u(0)\) ，并且 \(\varepsilon (t)\to 0\) 当 \(t\rightarrow +\infty\)。

\(f)\) 假设 \((\xi_j) \neq 0\) ；存在 \(f \in \mathrm{F}\) 使得积分 \(\int_{\mathrm{U}} f v_t\) 当 \(t \to +\infty\) 时不是有界的。（若该论断不真，则利用上一问题与巴拿赫–斯坦豪斯定理得出矛盾。）

\(g)\) 诸函数 \(v_{t}\) 在 \(\mathcal{D}'(\mathrm{U})\) 中不收敛于 0 当 \(t\to+\infty\) 时。\(h)\) 设 \(\xi=(0,0,\xi_{3})\in\mathbf{R}^{3}\) 满足 \(\xi_{3}<0\)。存在 \(\mathbf{R}^{3}\) 中 0 的邻域 U 以及函数 \(u\in\mathcal{C}^{\infty}(\mathrm{U})\) ，使得 \(\mathrm{L}(u)=0\) 且 \(u(0)=0\) ，并满足问题 \(b)\) 与 \(f)\) 的条件。（利用柯西–柯瓦列夫斯卡娅定理。）

\begin{enumerate}
\def\labelenumi{\roman{enumi})}
\tightlist
\item
  设 \((v_{t})\) 是 \(\mathcal{D}(\mathrm{U})\) 中满足 \(\mathrm{L}(v_{t})\to0\) 的函数，并设 \(f\in\mathcal{D}(\mathrm{U})\) 使得 \(\int_{\mathrm{U}}fv_{t}\) 不是有界的当 \(t\to+\infty\) 时（问题 g)）；方程 Lv=f 没有解 \(v\in\mathcal{D}'(\mathrm{U})\) （“卢伊定理”）。（注意 L 的形式伴随与 L 一致。）
\end{enumerate}

\begin{enumerate}
\def\labelenumi{\arabic{enumi})}
\setcounter{enumi}{31}
\tightlist
\item
  设 \(n \geqslant 1\) 是整数。我们用 \(\Delta\) 表示 \(\mathcal{D}'(\mathbf{R}^n)\) 上的拉普拉斯算子，它定义为
\end{enumerate}

\[
\Delta (f) = - \sum_ {i = 1} ^ {n} \partial_ {i} ^ {2} f.
\]

设 u 是 \(R^{n}\) 上的局部可积函数，使得 \(\Delta(u)\) 是局部可积函数。

设 \(v\in \mathcal{L}^{\infty}(\mathbf{R}^{n})\) 是由下式定义的函数：

\[
v (x) = \left\{ \begin{array}{l l} 0 & \text { if } u (x) = 0, \\ \overline {{u (x)}} / | u (x) | & \text { if } u (x) \neq 0. \end{array} \right.
\]

\begin{enumerate}
\def\labelenumi{\alph{enumi})}
\item
  对 \(\varepsilon > 0\)，令 \(u_{\varepsilon} = (|u|^2 + \varepsilon^2)^{1/2}\) 。证明广义函数 \(-\Delta(u_{\varepsilon}) - \mathcal{R}(\overline{u}/u_{\varepsilon} \Delta(u))\) 是正的。（先考虑 \(u \in \mathcal{C}^{\infty}(\mathbf{R}^{n})\) 的情形，再用卷积处理一般情形。）
\item
  广义函数 \(-\Delta(|u|) - \mathcal{R}(v\Delta(u))\) 是正的（“加藤不等式”）。
\end{enumerate}

\begin{enumerate}
\def\labelenumi{\arabic{enumi})}
\setcounter{enumi}{32}
\tightlist
\item
  设 \(m \geqslant 1\) 是整数。
\end{enumerate}

\begin{enumerate}
\def\labelenumi{\alph{enumi})}
\tightlist
\item
  存在唯一的线性映射 \(A^{+}\) （相应地 \(A^{-}\) ），它从 \(\mathrm{H}^{m}(\mathbf{R})\) 到 \(\mathrm{H}^{m-1}(\mathbf{R})\) ，使得
\end{enumerate}

\[
\mathrm{A} ^ {+} f (x) = f ^ {\prime} (x) - x f (x) \quad (\text { resp.   A } ^ {-} f (x) = f ^ {\prime} (x) + x f (x))
\]

对所有 \(f \in \mathcal{D}(\mathbf{R})\) 以及每个 \(x \in \mathbf{R}\) 成立。

\begin{enumerate}
\def\labelenumi{\alph{enumi})}
\setcounter{enumi}{1}
\tightlist
\item
  映射 \(A^{+}\) 是指标为 1 的弗雷德霍姆算子，而 \(A^{-}\) 是指标为 -1 的弗雷德霍姆算子。（考虑埃尔米特函数在 \(A^{+}\) 与 \(A^{-}\) 下的像，参见 II, p. 263 的习题 3。）
\end{enumerate}